% Fichier engendré par tools/generate-tex.py à partir de 11-informatique-nsi.
% Ne pas éditer à la main : tout le texte provient des commentaires du fichier Lean.
\documentclass[11pt,a4paper]{article}

% Compile aussi bien avec pdflatex qu'avec xelatex ou tectonic.
\usepackage{iftex}
\ifPDFTeX
  \usepackage[T1]{fontenc}
  \usepackage[utf8]{inputenc}
\else
  \usepackage{fontspec}
\fi
\usepackage[french]{babel}
\usepackage{amsmath,amssymb,amsthm}
\usepackage[margin=2.5cm]{geometry}
\usepackage[hidelinks]{hyperref}

% Les figures sont dessinées en TikZ : elles reprennent ainsi les
% polices du document. Voir la skill `illustrate-theorem`.
\usepackage{tikz}
\usetikzlibrary{calc}

\theoremstyle{definition}
\newtheorem{definition}{Définition}
\newtheorem{exemple}{Exemple}
\theoremstyle{plain}
\newtheorem{theoreme}{Théorème}
\newtheorem{lemme}[theoreme]{Lemme}

% Un exemple ne se démontre pas : on explique ce qu'il montre.
\newenvironment{explication}{\begin{proof}[Explication]}{\end{proof}}

% Un énoncé écrit mais non démontré : la preuve Lean est un `sorry`.
\newcommand{\admis}{\par\smallskip\noindent\textit{Admis : l'énoncé est écrit en Lean, sa démonstration reste à faire.}\par}

% Renvoi à la déclaration Lean, une ligne après la démonstration, aligné à droite.
\newcommand{\source}[2]{\par\smallskip\noindent\hfill{\footnotesize\href{#1}{\texttt{#2}}}\par\smallskip}

\title{Informatique}
\date{}

\begin{document}
\maketitle
% Les identifiants Lean en machine à écrire ne se coupent pas : on tolère des
% espacements plus lâches plutôt que des lignes qui débordent.
\sloppy

\section{RepresentationDesDonnees}

\noindent
Lycée \textemdash{} spécialité NSI, section \og{} Représentation des données \fg{}.
Un ordinateur ne manipule que des suites de bits ; tout le reste est convention.
Ce fichier démontre ces conventions : l'écriture binaire d'un entier existe et est
unique, elle se relit comme une somme de puissances de deux, un chiffre hexadécimal
vaut quatre bits, et un dixième n'a pas d'écriture binaire finie \textemdash{} d'où les surprises
du calcul flottant.

\subsection{Algèbre de Boole}

\begin{theoreme}
Pour deux booléens $a$ et $b$, la négation échange \og{} et \fg{} et \og{} ou \fg{} :
\[ \overline{a \wedge b} = \bar a \vee \bar b, \qquad \overline{a \vee b} = \bar a \wedge \bar b . \]
\end{theoreme}
\begin{proof}
Un booléen ne prend que deux valeurs : il y a donc quatre couples $(a, b)$ à examiner,
et chacun se vérifie par le calcul des deux membres. Pour $a$ vrai et $b$ faux, par
exemple, $a \wedge b$ est faux donc $\overline{a \wedge b}$ est vrai, tandis que
$\bar a$ est faux et $\bar b$ vrai, donc $\bar a \vee \bar b$ est vrai également.
\end{proof}
\source{https://github.com/Commutator-IO/learn-lean/blob/main/courses/02-lycee/11-informatique-nsi/RepresentationDesDonnees.lean\#L19}{RepresentationDesDonnees.lean\#L19}

\begin{theoreme}
Chaque opération se distribue sur l'autre, et la négation appliquée deux fois rend le
booléen de départ :
\[ a \wedge (b \vee c) = (a \wedge b) \vee (a \wedge c), \qquad
   a \vee (b \wedge c) = (a \vee b) \wedge (a \vee c), \qquad
   \bar{\bar a} = a . \]
\end{theoreme}
\begin{proof}
Même méthode : les trois booléens $a$, $b$, $c$ donnent huit cas, et chacun se vérifie
en calculant les deux membres. La seconde égalité mérite d'être remarquée : elle n'a pas
d'équivalent sur les nombres, où l'addition ne se distribue pas sur la multiplication.
\end{proof}
\source{https://github.com/Commutator-IO/learn-lean/blob/main/courses/02-lycee/11-informatique-nsi/RepresentationDesDonnees.lean\#L24}{RepresentationDesDonnees.lean\#L24}

\subsection{Écriture binaire}

\begin{theoreme}
Soit $n$ un entier naturel non nul. Sa suite de chiffres en base $2$ ne contient que
des $0$ et des $1$, et son dernier terme \textemdash{} le chiffre de poids fort \textemdash{}
n'est pas nul.
\end{theoreme}
\begin{proof}
Les deux points viennent de la construction des chiffres en base $b$ : chaque chiffre
est un reste de division par $b$, donc reste inférieur à $b$, ici à $2$ ; et le chiffre de
poids fort est non nul, faute de quoi le nombre s'écrirait avec un chiffre de moins.
C'est ce qui rend l'écriture canonique : on n'ajoute pas de zéros inutiles à gauche.
\end{proof}
\source{https://github.com/Commutator-IO/learn-lean/blob/main/courses/02-lycee/11-informatique-nsi/RepresentationDesDonnees.lean\#L33}{RepresentationDesDonnees.lean\#L33}

\begin{theoreme}
Relire une écriture binaire rend le nombre de départ :
\[ \sum_{i} b_i\, 2^i = n \quad \text{où } (b_i) \text{ est la suite des chiffres de } n . \]
\end{theoreme}
\begin{proof}
C'est la propriété qui justifie le mot \emph{écriture} : passer de $n$ à ses chiffres,
puis des chiffres à leur valeur, ramène à $n$. La démonstration est une récurrence sur le
nombre de chiffres, où le pas se lit sur la division euclidienne par $2$ : le dernier
chiffre est le reste, et les autres sont ceux du quotient.
\end{proof}
\source{https://github.com/Commutator-IO/learn-lean/blob/main/courses/02-lycee/11-informatique-nsi/RepresentationDesDonnees.lean\#L41}{RepresentationDesDonnees.lean\#L41}

\begin{theoreme}
L'écriture binaire est unique : si une suite $L$ ne contient que des $0$ et des $1$ et
ne se termine pas par un $0$, alors les chiffres de sa valeur sont exactement $L$.
\end{theoreme}
\begin{proof}
C'est la réciproque de l'énoncé précédent, et les deux hypothèses sont nécessaires. Sans
la première, une \og{} suite de chiffres \fg{} pourrait contenir un $2$ et désigner la
même valeur qu'une autre suite. Sans la seconde, $L$ et $L$ suivie d'un $0$ auraient la
même valeur : c'est le zéro inutile à gauche. Avec les deux, la division euclidienne par
$2$ redonne les chiffres un à un, dans l'ordre.
\end{proof}
\source{https://github.com/Commutator-IO/learn-lean/blob/main/courses/02-lycee/11-informatique-nsi/RepresentationDesDonnees.lean\#L47}{RepresentationDesDonnees.lean\#L47}

\begin{theoreme}
Un entier naturel non nul $n$ s'écrit avec $\lfloor \log_2 n \rfloor + 1$ bits.
\end{theoreme}
\begin{proof}
L'écriture de $n$ compte un chiffre par puissance de $2$ inférieure ou égale à $n$ :
elle en a $k+1$ exactement lorsque $2^k \le n < 2^{k+1}$, c'est-à-dire lorsque
$k = \lfloor \log_2 n \rfloor$. Un octet code donc les entiers de $0$ à $255$, et pas un
de plus.
\end{proof}
\source{https://github.com/Commutator-IO/learn-lean/blob/main/courses/02-lycee/11-informatique-nsi/RepresentationDesDonnees.lean\#L53}{RepresentationDesDonnees.lean\#L53}

\subsection{Complément à deux}

\begin{theoreme}
Sur $n$ bits, chaque entier relatif $a$ a un code et un seul parmi les $2^n$ codes
disponibles : l'unique entier $c$ tel que
\[ 0 \le c < 2^n \quad \text{et} \quad c \equiv a \pmod{2^n} , \]
et l'addition des codes se fait modulo $2^n$ :
\[ (a + b) \bmod 2^n = \bigl((a \bmod 2^n) + (b \bmod 2^n)\bigr) \bmod 2^n . \]
\end{theoreme}
\begin{proof}
Pour l'existence, $c = a \bmod 2^n$ convient : un reste de division par $2^n$ est compris
entre $0$ et $2^n - 1$, et il est congru à $a$ par définition. Pour l'unicité, si $c$
vérifie les deux conditions, alors $c$ étant déjà compris entre $0$ et $2^n$, son reste
modulo $2^n$ est lui-même ; la congruence donne donc $c = a \bmod 2^n$.

L'égalité sur l'addition est celle du reste d'une somme. C'est elle qui explique le
comportement des entiers machine : ajouter $1$ au plus grand entier positif d'un mot de
$n$ bits ne provoque pas d'erreur, mais rend le plus petit négatif \textemdash{} le calcul
est juste modulo $2^n$, et faux dans $\mathbb{Z}$.
\end{proof}
\source{https://github.com/Commutator-IO/learn-lean/blob/main/courses/02-lycee/11-informatique-nsi/RepresentationDesDonnees.lean\#L63}{RepresentationDesDonnees.lean\#L63}

\subsection{Hexadécimal}

\begin{theoreme}
Un chiffre hexadécimal vaut quatre bits : pour tous entiers $n$ et $k$,
\[ n < 16^k \iff n < 2^{4k} . \]
\end{theoreme}
\begin{proof}
Il suffit de remarquer que $16 = 2^4$, donc $16^k = (2^4)^k = 2^{4k}$ : les deux
inégalités portent sur le même nombre. C'est la raison pour laquelle l'hexadécimal sert à
noter les octets : deux chiffres hexadécimaux valent exactement huit bits.
\end{proof}
\source{https://github.com/Commutator-IO/learn-lean/blob/main/courses/02-lycee/11-informatique-nsi/RepresentationDesDonnees.lean\#L78}{RepresentationDesDonnees.lean\#L78}

\subsection{Nombres à virgule flottante}

\begin{theoreme}
Un dixième n'a pas d'écriture binaire finie : il n'existe pas d'entier relatif $a$ ni
d'entier naturel $n$ tels que
\[ \frac{1}{10} = \frac{a}{2^n} . \]
\end{theoreme}
\begin{proof}
Supposons une telle écriture. En multipliant les deux membres par $10 \times 2^n$, elle
donne $2^n = 10\,a$, donc $5$ divise $2^n$. Or $5$ est premier : s'il divise une puissance
de $2$, il divise $2$, ce qui est faux. L'hypothèse est donc absurde.

C'est l'origine des surprises du calcul flottant. La machine ne peut pas représenter
$0{,}1$ exactement ; elle en garde la valeur binaire la plus proche, et la somme
$0{,}1 + 0{,}2$ ne rend pas exactement $0{,}3$.
\end{proof}
\source{https://github.com/Commutator-IO/learn-lean/blob/main/courses/02-lycee/11-informatique-nsi/RepresentationDesDonnees.lean\#L89}{RepresentationDesDonnees.lean\#L89}

\subsection{Chaînes de caractères}

\begin{theoreme}
La concaténation des chaînes est associative, et la longueur d'une concaténation est la
somme des longueurs :
\[ (u + v) + w = u + (v + w), \qquad |u + v| = |u| + |v| . \]
\end{theoreme}
\begin{proof}
Les deux se démontrent par récurrence sur la première chaîne. Pour l'associativité, une
chaîne vide se retire des deux côtés, et le pas consiste à sortir le premier caractère,
qui se place en tête dans les deux membres. Pour la longueur, le même découpage ajoute
$1$ de chaque côté.
\end{proof}
\source{https://github.com/Commutator-IO/learn-lean/blob/main/courses/02-lycee/11-informatique-nsi/RepresentationDesDonnees.lean\#L108}{RepresentationDesDonnees.lean\#L108}

\section{AlgorithmesSurLesTableaux}

\noindent
Lycée \textemdash{} spécialité NSI, section \og{} Algorithmes sur les tableaux \fg{}.
Ce fichier écrit les algorithmes du programme \textemdash{} recherche séquentielle, recherche
dichotomique, tris par insertion, par sélection et par fusion \textemdash{} puis démontre pour
chacun les deux choses qu'on lui demande : qu'il rende le bon résultat, et combien il
coûte. Les coûts sont comptés par une fonction écrite à côté de l'algorithme, qui
compte les comparaisons ; c'est ce qui permet d'en parler sans arrondir.
Énoncés et démonstrations en français : voir Informatique.tex.

\subsection{Recherche séquentielle}

\begin{definition}
La recherche séquentielle : on parcourt le tableau du début à la fin et l'on rend
l'indice de la première case qui contient la valeur cherchée, ou rien.
\end{definition}
\source{https://github.com/Commutator-IO/learn-lean/blob/main/courses/02-lycee/11-informatique-nsi/AlgorithmesSurLesTableaux.lean\#L20}{AlgorithmesSurLesTableaux.lean\#L20}

\begin{theoreme}
La recherche séquentielle est correcte : quand elle rend un indice, la case
correspondante contient bien la valeur cherchée ; quand elle ne rend rien, la valeur
n'est pas dans le tableau. Les deux moitiés sont nécessaires \textemdash{} un programme qui
répond toujours \og{} absent \fg{} satisfait la première seule.
\end{theoreme}
\begin{proof}
Par récurrence sur le tableau.

S'il est vide, la recherche ne rend rien, et aucune valeur ne s'y trouve : les deux
moitiés sont vraies.

Sinon le tableau s'écrit \(a :: r\). Si \(a = x\), la recherche rend l'indice \(0\), et
la case \(0\) contient bien \(a\), c'est-à-dire \(x\).

Si \(a \ne x\), la recherche se poursuit sur \(r\) et ajoute \(1\) à ce qu'elle y
trouve. Si elle rend l'indice \(j+1\), c'est que la recherche dans \(r\) a rendu \(j\) ;
par hypothèse de récurrence la case \(j\) de \(r\) contient \(x\), et la case \(j+1\) de
\(a :: r\) est précisément la case \(j\) de \(r\). Si elle ne rend rien, c'est que la
recherche dans \(r\) n'a rien rendu ; par hypothèse de récurrence \(x\) n'est pas dans
\(r\), et comme \(a \ne x\), il n'est pas non plus dans \(a :: r\).
\end{proof}
\source{https://github.com/Commutator-IO/learn-lean/blob/main/courses/02-lycee/11-informatique-nsi/AlgorithmesSurLesTableaux.lean\#L28}{AlgorithmesSurLesTableaux.lean\#L28}

\begin{definition}
Le nombre de comparaisons que fait la recherche séquentielle : une par case
examinée, et l'on s'arrête dès qu'on a trouvé.
\end{definition}
\source{https://github.com/Commutator-IO/learn-lean/blob/main/courses/02-lycee/11-informatique-nsi/AlgorithmesSurLesTableaux.lean\#L53}{AlgorithmesSurLesTableaux.lean\#L53}

\begin{theoreme}
Le coût de la recherche séquentielle est d'au plus \texttt{n} comparaisons, et ce pire
cas est atteint : quand la valeur est absente, il faut examiner toutes les cases pour
pouvoir l'affirmer.
\end{theoreme}
\begin{proof}
Par récurrence sur le tableau. Sur le tableau vide, le coût et la longueur valent
\(0\).

Sur \(a :: r\), si \(a = x\), la recherche s'arrête après une comparaison, et la valeur
étant présente, la seconde moitié n'a rien à dire.

Si \(a \ne x\), le coût vaut \(1 + c(r)\), majoré par hypothèse de récurrence par
\(1 + \mathrm{longueur}(r)\), qui est la longueur de \(a :: r\). Si de plus \(x\) est
absent de \(a :: r\), il est absent de \(r\), donc \(c(r)\) vaut exactement la longueur
de \(r\) et le coût vaut celle du tableau.

Le pire cas est donc atteint, et il l'est précisément quand la valeur manque : c'est
l'absence qui coûte cher, puisqu'il faut avoir tout regardé pour l'affirmer.
\end{proof}
\source{https://github.com/Commutator-IO/learn-lean/blob/main/courses/02-lycee/11-informatique-nsi/AlgorithmesSurLesTableaux.lean\#L60}{AlgorithmesSurLesTableaux.lean\#L60}

\subsection{Maximum d'un tableau}

\begin{definition}
Le maximum, calculé en un parcours : on garde le plus grand élément vu jusque-là.
Le premier élément sert d'initialisation, ce qui impose que le tableau soit non vide
\textemdash{} un tableau vide n'a pas de maximum, et la signature le dit.
\end{definition}
\source{https://github.com/Commutator-IO/learn-lean/blob/main/courses/02-lycee/11-informatique-nsi/AlgorithmesSurLesTableaux.lean\#L78}{AlgorithmesSurLesTableaux.lean\#L78}

\begin{theoreme}
Le maximum d'un tableau non vide est bien un élément du tableau, et il majore
tous les autres. Là encore les deux moitiés sont nécessaires : sans la première, une
fonction qui rend un nombre très grand conviendrait.
\end{theoreme}
\begin{proof}
Par récurrence sur la queue du tableau, l'élément gardé restant libre — c'est
nécessaire, puisqu'il change à chaque tour.

Si la queue est vide, le maximum est l'élément gardé, seul élément du tableau.

Sinon le tableau s'écrit \(a :: b :: s\), et le calcul se poursuit avec \(\max(a, b)\)
gardé et \(s\) pour queue. L'hypothèse de récurrence donne alors deux choses : le
résultat \(M\) appartient à \(\max(a,b) :: s\), et il majore les éléments de cette
liste.

Or \(\max(a,b)\) vaut \(a\) ou \(b\) : dans les deux cas \(M\) appartient à
\(a :: b :: s\). Pour la majoration, \(a \le \max(a,b) \le M\) et
\(b \le \max(a,b) \le M\), et les éléments de \(s\) sont majorés par l'hypothèse de
récurrence.
\end{proof}
\source{https://github.com/Commutator-IO/learn-lean/blob/main/courses/02-lycee/11-informatique-nsi/AlgorithmesSurLesTableaux.lean\#L85}{AlgorithmesSurLesTableaux.lean\#L85}

\subsection{Recherche dichotomique}

\begin{definition}
La recherche dichotomique dans un tableau trié. Le tableau est ici donné par la
fonction \texttt{t} qui à un indice associe sa valeur, et l'on cherche dans la tranche
d'indices \texttt{[lo, hi[} : c'est exactement ce qu'on écrit en machine, où les deux bornes
sont deux variables qui se rapprochent.

On compare la valeur du milieu à celle qu'on cherche, et l'on ne garde que la moitié
où elle peut encore se trouver. La terminaison tient à ce que \texttt{hi − lo} diminue
strictement à chaque appel : c'est le variant de la boucle.
\end{definition}
\source{https://github.com/Commutator-IO/learn-lean/blob/main/courses/02-lycee/11-informatique-nsi/AlgorithmesSurLesTableaux.lean\#L116}{AlgorithmesSurLesTableaux.lean\#L116}

\begin{theoreme}
Quand la recherche dichotomique rend un indice, cette case contient bien la
valeur cherchée. Cette moitié-là ne demande pas que le tableau soit trié : on ne rend
un indice qu'après avoir comparé, et vu l'égalité.
\end{theoreme}
\begin{proof}
Par récurrence sur le carburant, c'est-à-dire sur un entier qui majore l'écart
\(hi - lo\) : c'est ainsi qu'on écrit une récurrence sur un variant qui n'est pas un
sous-terme.

Si \(lo \ge hi\), la tranche est vide et la recherche ne rend rien : il n'y a rien à
démontrer. Sinon on pose \(m = \lfloor (lo+hi)/2 \rfloor\). Si \(t(m) = x\), la
recherche rend \(m\), et la conclusion est exactement l'égalité qu'on vient de tester.
Sinon elle se poursuit sur l'une des deux moitiés, dont l'écart est strictement plus
petit, et l'hypothèse de récurrence s'applique.

Remarquons qu'à aucun moment on n'a utilisé que le tableau était trié.
\end{proof}
\source{https://github.com/Commutator-IO/learn-lean/blob/main/courses/02-lycee/11-informatique-nsi/AlgorithmesSurLesTableaux.lean\#L128}{AlgorithmesSurLesTableaux.lean\#L128}

\begin{theoreme}
Quand la recherche dichotomique ne rend rien, la valeur est absente de la tranche
examinée. C'est ici que le tri sert, et c'est tout ce qu'il sert : si la valeur du
milieu est trop petite, celles qui la précèdent le sont aussi et l'on peut jeter
toute la moitié gauche sans la regarder. Sur un tableau non trié, l'algorithme
répondrait \og{} absent \fg{} à tort.
\end{theoreme}
\begin{proof}
Par récurrence sur le carburant. Si \(lo \ge hi\), la tranche est vide : aucun indice
à considérer.

Sinon \(m = \lfloor (lo+hi)/2 \rfloor\) vérifie \(lo \le m < hi\). Comme la recherche
n'a rien rendu, elle n'a pas trouvé au milieu : \(t(m) \ne x\). Deux cas.

Si \(t(m) < x\), la recherche s'est poursuivie sur \([m+1, hi[\) sans rien rendre, et
l'hypothèse de récurrence en exclut \(x\). Reste la moitié gauche : pour
\(lo \le i \le m\), le tri donne \(t(i) \le t(m) < x\), donc \(t(i) \ne x\).

Si \(x < t(m)\), la recherche s'est poursuivie sur \([lo, m[\) ; l'hypothèse de
récurrence en exclut \(x\), et pour \(m \le i < hi\), le tri donne
\(x < t(m) \le t(i)\).

C'est le seul endroit de la démonstration où le tri intervient, et c'est bien pour cela
que la dichotomie ne s'applique qu'à un tableau trié : sur un tableau quelconque,
l'algorithme répondrait « absent » à tort.
\end{proof}
\source{https://github.com/Commutator-IO/learn-lean/blob/main/courses/02-lycee/11-informatique-nsi/AlgorithmesSurLesTableaux.lean\#L159}{AlgorithmesSurLesTableaux.lean\#L159}

\begin{definition}
Le nombre de comparaisons que fait la recherche dichotomique dans le pire des
cas : une par appel, et l'on continue dans la moitié la plus coûteuse.
\end{definition}
\source{https://github.com/Commutator-IO/learn-lean/blob/main/courses/02-lycee/11-informatique-nsi/AlgorithmesSurLesTableaux.lean\#L195}{AlgorithmesSurLesTableaux.lean\#L195}

\begin{theoreme}
La recherche dichotomique fait au plus \texttt{⌊log₂ n⌋ + 1} comparaisons sur une
tranche de \texttt{n} cases : chaque comparaison divise par deux le nombre de cases encore
possibles, et l'on ne peut diviser par deux qu'un nombre logarithmique de fois. C'est
tout l'intérêt du tri préalable : \texttt{20} comparaisons suffisent pour un million de
cases, contre un million pour la recherche séquentielle.
\end{theoreme}
\begin{proof}
Par récurrence sur le carburant. Sur une tranche vide, la recherche ne compare rien.

Sinon, posons \(n = hi - lo\) et \(m = \lfloor (lo+hi)/2 \rfloor\). Chacune des deux
moitiés compte au plus \(\lfloor n/2 \rfloor\) cases : à gauche
\(m - lo = \lfloor n/2 \rfloor\), à droite \(hi - m - 1 \le \lfloor n/2 \rfloor\).

Si \(n = 1\), le milieu est \(lo\), les deux moitiés sont vides, et le coût vaut \(1\),
qui est bien \(\log_2 1 + 1\).

Si \(n \ge 2\), l'hypothèse de récurrence majore le coût de chaque moitié par
\(\log_2 \lfloor n/2 \rfloor + 1\), et la croissance du logarithme donne
\(\log_2 \lfloor n/2 \rfloor \le \log_2 n - 1\). En ajoutant la comparaison du tour
courant :
\[ 1 + \bigl(\log_2 n - 1 + 1\bigr) = \log_2 n + 1 . \]

C'est tout l'intérêt du tri préalable : vingt comparaisons suffisent pour un million de
cases, là où la recherche séquentielle en demanderait un million.
\end{proof}
\source{https://github.com/Commutator-IO/learn-lean/blob/main/courses/02-lycee/11-informatique-nsi/AlgorithmesSurLesTableaux.lean\#L207}{AlgorithmesSurLesTableaux.lean\#L207}

\subsection{Tri par insertion}

\begin{definition}
L'insertion d'un élément dans une liste déjà triée : on descend jusqu'à la
première place où il n'est pas trop grand, et on l'y pose.
\end{definition}
\source{https://github.com/Commutator-IO/learn-lean/blob/main/courses/02-lycee/11-informatique-nsi/AlgorithmesSurLesTableaux.lean\#L248}{AlgorithmesSurLesTableaux.lean\#L248}

\begin{definition}
Le tri par insertion : on trie la fin du tableau, puis on insère le premier
élément à sa place.
\end{definition}
\source{https://github.com/Commutator-IO/learn-lean/blob/main/courses/02-lycee/11-informatique-nsi/AlgorithmesSurLesTableaux.lean\#L254}{AlgorithmesSurLesTableaux.lean\#L254}

\begin{theoreme}
L'insertion ne perd ni n'invente d'élément : la liste obtenue est une permutation
de la liste de départ augmentée de l'élément inséré.
\end{theoreme}
\begin{proof}
Par récurrence sur la liste. Insérer dans la liste vide donne \([x]\).

Sur \(a :: r\), si \(x \le a\), l'insertion pose \(x\) en tête et ne touche à rien
d'autre. Sinon elle rend \(a :: \mathrm{insere}(x, r)\), qui par hypothèse de
récurrence est une permutation de \(a :: x :: r\) ; échanger les deux premiers éléments
donne \(x :: a :: r\).
\end{proof}
\source{https://github.com/Commutator-IO/learn-lean/blob/main/courses/02-lycee/11-informatique-nsi/AlgorithmesSurLesTableaux.lean\#L260}{AlgorithmesSurLesTableaux.lean\#L260}

\begin{theoreme}
Les éléments de la liste après insertion sont ceux d'avant, plus l'inséré.
\end{theoreme}
\begin{proof}
Deux listes permutées ont les mêmes éléments : c'est l'énoncé précédent, lu au niveau
des appartenances.
\end{proof}
\source{https://github.com/Commutator-IO/learn-lean/blob/main/courses/02-lycee/11-informatique-nsi/AlgorithmesSurLesTableaux.lean\#L271}{AlgorithmesSurLesTableaux.lean\#L271}

\begin{theoreme}
L'insertion dans une liste triée rend une liste triée : c'est l'unique chose à
vérifier, puisque le tri se contente de l'appliquer élément par élément.
\end{theoreme}
\begin{proof}
Par récurrence sur la liste. Dans la liste vide, \([x]\) est trié.

Sinon la liste s'écrit \(a :: r\), où \(a\) minore les éléments de \(r\) et où \(r\) est
triée.

Si \(x \le a\), le résultat est \(x :: a :: r\), qui est trié : \(x \le a\), et pour
tout \(y\) de \(r\), \(x \le a \le y\).

Sinon \(a < x\), et le résultat est \(a :: \mathrm{insere}(x, r)\). Par hypothèse de
récurrence \(\mathrm{insere}(x, r)\) est triée, et ses éléments sont ceux de \(r\)
augmentés de \(x\) : tous sont supérieurs à \(a\) — ceux de \(r\) parce que la liste de
départ était triée, et \(x\) parce que \(a < x\).
\end{proof}
\source{https://github.com/Commutator-IO/learn-lean/blob/main/courses/02-lycee/11-informatique-nsi/AlgorithmesSurLesTableaux.lean\#L278}{AlgorithmesSurLesTableaux.lean\#L278}

\begin{theoreme}
Le tri par insertion est correct : son résultat est trié, et c'est une
permutation de l'entrée. Les deux moitiés sont nécessaires \textemdash{} la liste vide est triée,
et la liste de départ est une permutation d'elle-même.
\end{theoreme}
\begin{proof}
Par récurrence sur la liste. La liste vide est triée et se permute en elle-même.

Sinon le tri de \(a :: r\) est l'insertion de \(a\) dans le tri de \(r\). Par hypothèse
de récurrence ce dernier est trié, et l'insertion dans une liste triée rend une liste
triée. Enfin l'insertion rend une permutation de \(a :: \mathrm{triInsertion}(r)\),
elle-même permutation de \(a :: r\) puisque le tri de \(r\) permute \(r\).
\end{proof}
\source{https://github.com/Commutator-IO/learn-lean/blob/main/courses/02-lycee/11-informatique-nsi/AlgorithmesSurLesTableaux.lean\#L302}{AlgorithmesSurLesTableaux.lean\#L302}

\subsection{Tri par sélection}

\begin{definition}
L'extraction du minimum : on parcourt la liste en gardant le plus petit élément
vu, et l'on rend ce minimum accompagné de tous les autres éléments.
\end{definition}
\source{https://github.com/Commutator-IO/learn-lean/blob/main/courses/02-lycee/11-informatique-nsi/AlgorithmesSurLesTableaux.lean\#L316}{AlgorithmesSurLesTableaux.lean\#L316}

\begin{theoreme}
L'extraction rend autant d'éléments qu'elle en a reçu, moins le minimum : c'est
ce qui fait que le tri par sélection termine.
\end{theoreme}
\begin{proof}
Par récurrence sur la liste, l'élément gardé restant libre. Sur la liste vide, le reste
est vide.

Sur \(a :: r\), quelle que soit la branche prise, le reste rendu est un élément suivi du
reste extrait de \(r\) : sa longueur vaut \(1 + \mathrm{longueur}(r)\), c'est-à-dire la
longueur de \(a :: r\).

C'est cette égalité qui justifie la terminaison du tri par sélection : sans elle, rien
ne dirait que la liste sur laquelle on recommence est plus courte.
\end{proof}
\source{https://github.com/Commutator-IO/learn-lean/blob/main/courses/02-lycee/11-informatique-nsi/AlgorithmesSurLesTableaux.lean\#L324}{AlgorithmesSurLesTableaux.lean\#L324}

\begin{definition}
Le tri par sélection : on extrait le minimum, on le pose en tête, et l'on
recommence sur le reste. La récursion ne porte pas sur un sous-terme mais sur une
liste de longueur strictement plus petite : il faut le dire à Lean.
\end{definition}
\source{https://github.com/Commutator-IO/learn-lean/blob/main/courses/02-lycee/11-informatique-nsi/AlgorithmesSurLesTableaux.lean\#L334}{AlgorithmesSurLesTableaux.lean\#L334}

\begin{theoreme}
L'extraction ne perd ni n'invente d'élément.
\end{theoreme}
\begin{proof}
Par récurrence sur la liste, l'élément gardé restant libre. Sur la liste vide, il n'y a
rien à faire.

Sur \(a :: r\), deux cas. Si \(a < m\), l'extraction se poursuit avec \(a\) gardé et
rend le minimum \(p\) accompagné de \(m\) suivi du reste \(s\). Par hypothèse de
récurrence \(p :: s\) permute \(a :: r\), donc
\[ p :: m :: s \sim m :: p :: s \sim m :: a :: r , \]
la première étape échangeant les deux premiers éléments.

Si \(m \le a\), l'extraction se poursuit avec \(m\) gardé et rend \(p\) accompagné de
\(a\) suivi du reste ; le même raisonnement s'applique, avec un échange de plus pour
remettre \(m\) en tête.
\end{proof}
\source{https://github.com/Commutator-IO/learn-lean/blob/main/courses/02-lycee/11-informatique-nsi/AlgorithmesSurLesTableaux.lean\#L341}{AlgorithmesSurLesTableaux.lean\#L341}

\begin{theoreme}
Le minimum extrait minore effectivement tous les éléments reçus.
\end{theoreme}
\begin{proof}
Par récurrence sur la liste, l'élément gardé restant libre. Sur la liste vide, le
minimum est l'élément gardé, qui se minore lui-même.

Sur \(a :: r\), deux cas. Si \(a < m\), le minimum rendu est celui extrait de \(a\) et
\(r\) : par hypothèse de récurrence il minore \(a\) et les éléments de \(r\). Il minore
aussi \(m\), puisqu'il est inférieur à \(a\), lui-même inférieur à \(m\).

Si \(m \le a\), le minimum rendu est celui extrait de \(m\) et \(r\) : il minore \(m\)
et les éléments de \(r\), donc aussi \(a\), qui majore \(m\).
\end{proof}
\source{https://github.com/Commutator-IO/learn-lean/blob/main/courses/02-lycee/11-informatique-nsi/AlgorithmesSurLesTableaux.lean\#L353}{AlgorithmesSurLesTableaux.lean\#L353}

\begin{theoreme}
Le tri par sélection est correct : même énoncé que pour le tri par insertion, et
c'est le point de l'exercice \textemdash{} deux programmes de structures très différentes
satisfont la même spécification.
\end{theoreme}
\begin{proof}
Par récurrence sur la longueur de la liste — la récursion du tri par sélection ne porte
pas sur un sous-terme, mais sur une liste plus courte.

La liste vide est triée et se permute en elle-même. Sinon on extrait le minimum \(p\) et
le reste \(s\), et le résultat est \(p :: \mathrm{triSelection}(s)\). Par hypothèse de
récurrence, \(\mathrm{triSelection}(s)\) est triée et permute \(s\).

Reste à voir que \(p\) minore ses éléments. Un élément de \(\mathrm{triSelection}(s)\)
est un élément de \(s\), donc de \(p :: s\), qui permute la liste de départ ; or \(p\)
minore tous les éléments de celle-ci.

Enfin \(p :: \mathrm{triSelection}(s) \sim p :: s \sim a :: r\).

C'est le point de l'exercice : deux programmes de structures très différentes — l'un
insère, l'autre sélectionne — satisfont la même spécification.
\end{proof}
\source{https://github.com/Commutator-IO/learn-lean/blob/main/courses/02-lycee/11-informatique-nsi/AlgorithmesSurLesTableaux.lean\#L377}{AlgorithmesSurLesTableaux.lean\#L377}

\subsection{Coût des tris quadratiques}

\begin{definition}
Le nombre de comparaisons que fait une insertion.
\end{definition}
\source{https://github.com/Commutator-IO/learn-lean/blob/main/courses/02-lycee/11-informatique-nsi/AlgorithmesSurLesTableaux.lean\#L394}{AlgorithmesSurLesTableaux.lean\#L394}

\begin{definition}
Le nombre de comparaisons que fait le tri par insertion.
\end{definition}
\source{https://github.com/Commutator-IO/learn-lean/blob/main/courses/02-lycee/11-informatique-nsi/AlgorithmesSurLesTableaux.lean\#L399}{AlgorithmesSurLesTableaux.lean\#L399}

\begin{theoreme}
Une insertion coûte au plus la longueur de la liste où l'on insère.
\end{theoreme}
\begin{proof}
Par récurrence : une comparaison par élément traversé, et l'insertion s'arrête au plus
tard à la fin de la liste.
\end{proof}
\source{https://github.com/Commutator-IO/learn-lean/blob/main/courses/02-lycee/11-informatique-nsi/AlgorithmesSurLesTableaux.lean\#L404}{AlgorithmesSurLesTableaux.lean\#L404}

\begin{theoreme}
Le tri par insertion fait au plus \texttt{n(n−1)/2} comparaisons. L'inégalité est écrite
sous la forme \texttt{2c + n \(\le\) n²} pour éviter à la fois la division et la soustraction
tronquée des entiers naturels.
\end{theoreme}
\begin{proof}
Par récurrence sur la liste. Sur la liste vide, tout est nul.

Sur \(a :: r\), notons \(k\) la longueur de \(r\). Le coût est celui du tri de \(r\),
augmenté de celui de l'insertion de \(a\) dans \(\mathrm{triInsertion}(r)\). Cette liste
a \(k\) éléments — trier n'en change pas le nombre —, donc l'insertion coûte au plus
\(k\). En notant \(c\) le coût du tri de \(r\), l'hypothèse de récurrence donne
\(2c + k \le k^2\), d'où
\[ 2\,c(a :: r) + (k+1) \le 2c + 2k + k + 1 \le k^2 + 2k + 1 = (k+1)^2 . \]

L'inégalité est écrite sous la forme \(2c + n \le n^2\) pour éviter à la fois la
division et la soustraction tronquée des entiers naturels ; elle dit exactement
\(c \le n(n-1)/2\).
\end{proof}
\source{https://github.com/Commutator-IO/learn-lean/blob/main/courses/02-lycee/11-informatique-nsi/AlgorithmesSurLesTableaux.lean\#L417}{AlgorithmesSurLesTableaux.lean\#L417}

\begin{definition}
Une liste strictement décroissante : \texttt{n, n−1, \(\dots\), 1}. C'est le pire cas du tri par
insertion, puisque chaque élément inséré est plus grand que tous ceux déjà triés et
doit donc les traverser tous.
\end{definition}
\source{https://github.com/Commutator-IO/learn-lean/blob/main/courses/02-lycee/11-informatique-nsi/AlgorithmesSurLesTableaux.lean\#L434}{AlgorithmesSurLesTableaux.lean\#L434}

\begin{theoreme}
Une insertion parcourt toute la liste quand l'élément inséré est plus grand que
tous les autres.
\end{theoreme}
\begin{proof}
Par récurrence. Si l'élément inséré est strictement plus grand que tous ceux de la
liste, aucune comparaison ne réussit, et l'insertion traverse la liste entière.
\end{proof}
\source{https://github.com/Commutator-IO/learn-lean/blob/main/courses/02-lycee/11-informatique-nsi/AlgorithmesSurLesTableaux.lean\#L440}{AlgorithmesSurLesTableaux.lean\#L440}

\begin{theoreme}
Les éléments de \texttt{listeDecroissante n} sont les entiers de \texttt{1} à \texttt{n}.
\end{theoreme}
\begin{proof}
Par récurrence : la liste de rang \(k+1\) ajoute l'élément \(k+1\) à celle de rang
\(k\), dont les éléments sont compris entre \(1\) et \(k\).
\end{proof}
\source{https://github.com/Commutator-IO/learn-lean/blob/main/courses/02-lycee/11-informatique-nsi/AlgorithmesSurLesTableaux.lean\#L450}{AlgorithmesSurLesTableaux.lean\#L450}

\begin{theoreme}
La longueur de \texttt{listeDecroissante n} est \texttt{n}.
\end{theoreme}
\begin{proof}
Par récurrence : un élément de plus à chaque rang.
\end{proof}
\source{https://github.com/Commutator-IO/learn-lean/blob/main/courses/02-lycee/11-informatique-nsi/AlgorithmesSurLesTableaux.lean\#L459}{AlgorithmesSurLesTableaux.lean\#L459}

\begin{theoreme}
Le pire cas quadratique est atteint : sur une liste strictement décroissante, le
tri par insertion fait exactement \texttt{n(n−1)/2} comparaisons. La majoration précédente
n'est donc pas améliorable.
\end{theoreme}
\begin{proof}
Par récurrence. Au rang \(0\), tout est nul.

Au rang \(k+1\), la liste est \((k+1) :: \mathrm{listeDecroissante}(k)\). Le tri de la
queue a \(k\) éléments, tous compris entre \(1\) et \(k\), donc tous strictement
inférieurs à \(k+1\) : l'insertion de \(k+1\) les traverse tous et coûte exactement
\(k\). D'où, en notant \(c\) le coût au rang \(k\) :
\[ 2\,c(k+1) + (k+1) = 2c + 2k + k + 1 = k^2 + 2k + 1 = (k+1)^2 . \]

La majoration précédente n'est donc pas améliorable : le pire cas quadratique est
atteint, et il l'est sur un tableau rangé à l'envers.
\end{proof}
\source{https://github.com/Commutator-IO/learn-lean/blob/main/courses/02-lycee/11-informatique-nsi/AlgorithmesSurLesTableaux.lean\#L467}{AlgorithmesSurLesTableaux.lean\#L467}

\subsection{Tri fusion}

\begin{definition}
La fusion de deux listes triées : on compare les deux têtes et l'on prend la plus
petite. La récursion porte sur la somme des deux longueurs, qui diminue à chaque
appel.
\end{definition}
\source{https://github.com/Commutator-IO/learn-lean/blob/main/courses/02-lycee/11-informatique-nsi/AlgorithmesSurLesTableaux.lean\#L490}{AlgorithmesSurLesTableaux.lean\#L490}

\begin{theoreme}
La fusion ne perd ni n'invente d'élément.
\end{theoreme}
\begin{proof}
Par récurrence sur la somme des deux longueurs, qui diminue à chaque appel.

Si l'une des deux listes est vide, la fusion rend l'autre. Sinon elle prend la plus
petite des deux têtes. Si c'est \(a\), le reste de la fusion permute, par hypothèse de
récurrence, la concaténation de \(r\) et de \(b :: s\), et il suffit de remettre \(a\)
devant. Si c'est \(b\), le reste permute la concaténation de \(a :: r\) et de \(s\), et
remettre \(b\) devant demande de le faire traverser la première liste — ce qui est bien
une permutation.
\end{proof}
\source{https://github.com/Commutator-IO/learn-lean/blob/main/courses/02-lycee/11-informatique-nsi/AlgorithmesSurLesTableaux.lean\#L497}{AlgorithmesSurLesTableaux.lean\#L497}

\begin{theoreme}
La fusion de deux listes triées est triée.
\end{theoreme}
\begin{proof}
Par récurrence sur la somme des deux longueurs. Si l'une est vide, le résultat est
l'autre, triée par hypothèse.

Sinon, disons \(a \le b\) : le résultat est \(a :: \mathrm{fusion}(r, b :: s)\), trié
par hypothèse de récurrence. Reste à voir que \(a\) minore ses éléments. Ce sont ceux de
\(r\) et ceux de \(b :: s\), puisque la fusion ne perd ni n'invente rien : ceux de \(r\)
sont majorés par \(a\) parce que \(a :: r\) est triée, et ceux de \(b :: s\) sont
supérieurs à \(b\), lui-même supérieur à \(a\).

Le cas \(b < a\) est symétrique.
\end{proof}
\source{https://github.com/Commutator-IO/learn-lean/blob/main/courses/02-lycee/11-informatique-nsi/AlgorithmesSurLesTableaux.lean\#L511}{AlgorithmesSurLesTableaux.lean\#L511}

\begin{definition}
Le tri fusion : on coupe la liste en deux moitiés, on trie chacune, on fusionne.
La récursion porte sur la longueur ; couper au milieu ne suffit pas à la faire
diminuer si la liste a moins de deux éléments, d'où les deux cas de base.
\end{definition}
\source{https://github.com/Commutator-IO/learn-lean/blob/main/courses/02-lycee/11-informatique-nsi/AlgorithmesSurLesTableaux.lean\#L543}{AlgorithmesSurLesTableaux.lean\#L543}

\begin{theoreme}
Le tri fusion est correct : son résultat est trié, et c'est une permutation de
l'entrée. La démonstration se ramène aux deux propriétés de la fusion, et à ce que
recoller les deux moitiés redonne la liste de départ.
\end{theoreme}
\begin{proof}
Par récurrence sur la longueur. Les listes de zéro ou un élément sont déjà triées.

Sinon on coupe la liste en deux moitiés, chacune strictement plus courte que la liste
entière — c'est ce qui fait terminer l'algorithme, et c'est pourquoi il faut deux cas de
base : couper au milieu une liste d'un seul élément ne la raccourcirait pas. Par
hypothèse de récurrence, le tri de chaque moitié est trié et la permute.

La fusion de deux listes triées est triée, et elle permute leur concaténation. Or
recoller les deux moitiés d'une liste redonne la liste.
\end{proof}
\source{https://github.com/Commutator-IO/learn-lean/blob/main/courses/02-lycee/11-informatique-nsi/AlgorithmesSurLesTableaux.lean\#L557}{AlgorithmesSurLesTableaux.lean\#L557}

\begin{theoreme}
Trier ne change pas le nombre d'éléments.
\end{theoreme}
\begin{proof}
Deux listes permutées ont même longueur.
\end{proof}
\source{https://github.com/Commutator-IO/learn-lean/blob/main/courses/02-lycee/11-informatique-nsi/AlgorithmesSurLesTableaux.lean\#L572}{AlgorithmesSurLesTableaux.lean\#L572}

\begin{definition}
Le nombre de comparaisons que fait la fusion : une par tour, et l'on s'arrête dès
qu'une des deux listes est épuisée \textemdash{} la fin de l'autre est recopiée sans comparaison.
\end{definition}
\source{https://github.com/Commutator-IO/learn-lean/blob/main/courses/02-lycee/11-informatique-nsi/AlgorithmesSurLesTableaux.lean\#L578}{AlgorithmesSurLesTableaux.lean\#L578}

\begin{theoreme}
La fusion est linéaire : elle fait au plus autant de comparaisons que la somme des
longueurs. C'est ce qui rend le tri fusion intéressant \textemdash{} le recollement ne coûte pas
plus cher qu'un parcours.
\end{theoreme}
\begin{proof}
Par récurrence sur la somme des deux longueurs. Si l'une est vide, la fusion ne compare
rien : la fin de l'autre est recopiée telle quelle. Sinon chaque tour consomme une
comparaison et une tête, quelle que soit la branche prise.

C'est ce qui rend le tri fusion intéressant : le recollement ne coûte pas plus cher
qu'un parcours.
\end{proof}
\source{https://github.com/Commutator-IO/learn-lean/blob/main/courses/02-lycee/11-informatique-nsi/AlgorithmesSurLesTableaux.lean\#L587}{AlgorithmesSurLesTableaux.lean\#L587}

\begin{definition}
Le nombre de comparaisons que fait le tri fusion : celles des deux moitiés, plus
celles du recollement.
\end{definition}
\source{https://github.com/Commutator-IO/learn-lean/blob/main/courses/02-lycee/11-informatique-nsi/AlgorithmesSurLesTableaux.lean\#L600}{AlgorithmesSurLesTableaux.lean\#L600}

\begin{theoreme}
Le tri fusion coûte \texttt{n log₂ n} comparaisons : au plus \texttt{n ⌈log₂ n⌉}.

C'est la résolution de \texttt{T(n) = 2T(n/2) + n} faite sur le programme lui-même, et non
sur la seule récurrence : il y a \texttt{⌈log₂ n⌉} étages de découpe, et les recollements
d'un même étage portent ensemble sur les \texttt{n} éléments, d'où \texttt{n} par étage.

Le logarithme est ici arrondi par excès, ce qui est nécessaire : pour \texttt{n = 3}, on
découpe en \texttt{1} et \texttt{2}, et il faut bien deux étages.
\end{theoreme}
\begin{proof}
Par récurrence sur la longueur. Les listes de zéro ou un élément ne coûtent rien.

Soit \(n \ge 2\) la longueur. Les deux moitiés comptent \(\lfloor n/2 \rfloor\) et
\(\lceil n/2 \rceil\) éléments, et le recollement coûte au plus \(n\) comparaisons
puisque la fusion est linéaire. Posons \(c = \lceil \log_2 n \rceil\) ; alors
\(c \ge 1\), et
\[ \bigl\lceil \log_2 \lceil n/2 \rceil \bigr\rceil = c - 1 , \qquad
   \bigl\lceil \log_2 \lfloor n/2 \rfloor \bigr\rceil \le c - 1 , \]
la seconde inégalité par croissance du logarithme. L'hypothèse de récurrence donne
alors
\[ \lfloor n/2 \rfloor (c-1) + \lceil n/2 \rceil (c-1) + n = n(c-1) + n = nc . \]

Le logarithme est ici arrondi par excès, et il le faut : pour \(n = 3\), on découpe en
\(1\) et \(2\), ce qui demande bien deux étages.

C'est la récurrence \(T(n) = 2T(n/2) + n\) résolue sur le programme lui-même, et non
sur la seule formule : il y a \(\lceil \log_2 n \rceil\) étages de découpe, et les
recollements d'un même étage portent ensemble sur les \(n\) éléments.
\end{proof}
\source{https://github.com/Commutator-IO/learn-lean/blob/main/courses/02-lycee/11-informatique-nsi/AlgorithmesSurLesTableaux.lean\#L621}{AlgorithmesSurLesTableaux.lean\#L621}

\subsection{Borne inférieure des tris par comparaisons}

\begin{definition}
L'arbre des comparaisons d'un algorithme de tri : chaque nœud est une comparaison
et ses deux enfants sont les deux réponses possibles ; chaque feuille est une sortie
de l'algorithme, donc une permutation qu'il peut produire.
\end{definition}
\source{https://github.com/Commutator-IO/learn-lean/blob/main/courses/02-lycee/11-informatique-nsi/AlgorithmesSurLesTableaux.lean\#L681}{AlgorithmesSurLesTableaux.lean\#L681}

\begin{definition}
La hauteur de l'arbre, c'est-à-dire le nombre de comparaisons du pire cas.
\end{definition}
\source{https://github.com/Commutator-IO/learn-lean/blob/main/courses/02-lycee/11-informatique-nsi/AlgorithmesSurLesTableaux.lean\#L686}{AlgorithmesSurLesTableaux.lean\#L686}

\begin{definition}
Le nombre de feuilles, c'est-à-dire le nombre de résultats distincts possibles.
\end{definition}
\source{https://github.com/Commutator-IO/learn-lean/blob/main/courses/02-lycee/11-informatique-nsi/AlgorithmesSurLesTableaux.lean\#L691}{AlgorithmesSurLesTableaux.lean\#L691}

\begin{theoreme}
Un arbre de hauteur \texttt{h} a au plus \texttt{2ʰ} feuilles : chaque comparaison ne fait au
mieux que doubler le nombre de cas qu'on sait distinguer.
\end{theoreme}
\begin{proof}
Par récurrence sur l'arbre. Une feuille a une feuille et la hauteur \(0\), et
\(1 \le 2^0\).

Un nœud de sous-arbres \(g\) et \(d\) a \(f(g) + f(d)\) feuilles et la hauteur
\(1 + \max(h(g), h(d))\). Par hypothèse de récurrence,
\[ f(g) + f(d) \le 2^{h(g)} + 2^{h(d)} \le 2 \times 2^{\max(h(g),\, h(d))}
   = 2^{1 + \max(h(g),\, h(d))} . \]

C'est la traduction du fait qu'une comparaison ne fait, au mieux, que doubler le nombre
de cas qu'on sait distinguer.
\end{proof}
\source{https://github.com/Commutator-IO/learn-lean/blob/main/courses/02-lycee/11-informatique-nsi/AlgorithmesSurLesTableaux.lean\#L697}{AlgorithmesSurLesTableaux.lean\#L697}

\begin{theoreme}
Un tri par comparaisons demande au moins \texttt{log₂(n!)} comparaisons.

L'hypothèse \texttt{hsorties} est ce qu'on demande à l'algorithme : pour trier
correctement \texttt{n} éléments, il doit pouvoir produire chacune des \texttt{n!} permutations,
donc son arbre a au moins \texttt{n!} feuilles. Comme un arbre de hauteur \texttt{h} n'a jamais
plus de \texttt{2ʰ} feuilles, il vient \texttt{n! \(\le\) 2ʰ}, puis la minoration par le logarithme.

Ce que cet énoncé ne formalise pas : que tout algorithme de tri par comparaisons
s'écrive bien sous la forme d'un tel arbre. C'est une modélisation, et elle est ici
portée par l'hypothèse plutôt que passée sous silence.
\end{theoreme}
\begin{proof}
L'arbre de l'algorithme a au moins \(n!\) feuilles par hypothèse, et au plus \(2^h\) par
ce qui précède, où \(h\) est sa hauteur. Donc \(n! \le 2^h\), et la croissance du
logarithme donne
\[ \log_2(n!) \le \log_2\bigl(2^h\bigr) = h . \]

C'est une minoration du pire cas, non du cas moyen : elle dit qu'aucun tri par
comparaisons ne peut faire mieux que \(\log_2(n!)\) comparaisons sur la plus mauvaise de
ses entrées. Comme \(\log_2(n!)\) vaut de l'ordre de \(n \log_2 n\), le tri fusion est
optimal à une constante près.

Ce que cet énoncé ne formalise pas : que tout algorithme de tri par comparaisons
s'écrive sous la forme d'un tel arbre. C'est une modélisation, portée ici par
l'hypothèse plutôt que passée sous silence.
\end{proof}
\source{https://github.com/Commutator-IO/learn-lean/blob/main/courses/02-lycee/11-informatique-nsi/AlgorithmesSurLesTableaux.lean\#L718}{AlgorithmesSurLesTableaux.lean\#L718}

\section{CorrectionEtTerminaisonDesProgrammes}

\noindent
Lycée \textemdash{} spécialité NSI, section \og{} Correction et terminaison des programmes \fg{}.
Un programme n'est pas correct parce qu'il donne le bon résultat sur les exemples
essayés : il l'est parce qu'on sait dire ce qu'il conserve à chaque tour de boucle,
et pourquoi il s'arrête. Ce fichier démontre les deux outils de ce raisonnement \textemdash{}
l'invariant et le variant \textemdash{} puis s'en sert sur deux programmes du programme de
première, avant de montrer que ce raisonnement a une limite : aucun programme ne
peut décider si un autre s'arrête.
Énoncés et démonstrations en français : voir Informatique.tex.

\subsection{Invariant de boucle}

\begin{theoreme}
Une boucle est modélisée par le corps qu'elle répète : \texttt{corps\textasciicircum{}[n] e} est l'état
obtenu après \texttt{n} tours à partir de l'état \texttt{e}.

Invariant de boucle : une propriété vraie avant la boucle et conservée par le corps
est vraie après un nombre quelconque de tours, donc en particulier à la sortie.
\end{theoreme}
\begin{proof}
Par récurrence sur le nombre de tours. Au tour \(0\), l'état est celui de départ, et la
propriété y est vraie par hypothèse. Si elle est vraie après \(k\) tours, elle l'est
après \(k+1\) puisque le corps la conserve.

C'est le raisonnement qu'on mène à la main sur un programme, écrit une fois pour
toutes : il n'y a rien de plus dans un invariant de boucle.
\end{proof}
\source{https://github.com/Commutator-IO/learn-lean/blob/main/courses/02-lycee/11-informatique-nsi/CorrectionEtTerminaisonDesProgrammes.lean\#L22}{CorrectionEtTerminaisonDesProgrammes.lean\#L22}

\subsection{Variant de boucle}

\begin{theoreme}
Variant de boucle : si l'on sait associer à l'état un entier naturel qui décroît
strictement à chaque tour tant que la condition est vraie, alors la condition finit
par devenir fausse \textemdash{} la boucle s'arrête.

C'est le seul endroit où l'on se sert de ce que \(\mathbb{N}\) n'a pas de suite infinie
strictement décroissante : un entier naturel qui diminue à chaque tour ne peut pas
diminuer indéfiniment.
\end{theoreme}
\begin{proof}
Par récurrence forte sur la valeur du variant à l'entrée.

Si la condition est déjà fausse, on s'arrête au tour \(0\). Sinon on fait un tour : le
variant a strictement diminué, donc l'hypothèse de récurrence s'applique à l'état
obtenu et fournit un nombre de tours après lequel la condition est fausse ; il suffit
d'en ajouter un.

Tout repose sur ce que les entiers naturels n'admettent pas de suite infinie
strictement décroissante. Un variant à valeurs réelles, ou entières relatives, ne
prouverait rien : on peut décroître indéfiniment.
\end{proof}
\source{https://github.com/Commutator-IO/learn-lean/blob/main/courses/02-lycee/11-informatique-nsi/CorrectionEtTerminaisonDesProgrammes.lean\#L40}{CorrectionEtTerminaisonDesProgrammes.lean\#L40}

\subsection{Somme des premiers entiers par accumulation}

\begin{definition}
Le corps de la boucle qui somme les entiers un à un. L'état est le couple
\texttt{(i, s)} : \texttt{i} compte les tours, \texttt{s} accumule la somme. Un tour ajoute \texttt{i + 1} à
l'accumulateur et avance le compteur.
\end{definition}
\source{https://github.com/Commutator-IO/learn-lean/blob/main/courses/02-lycee/11-informatique-nsi/CorrectionEtTerminaisonDesProgrammes.lean\#L57}{CorrectionEtTerminaisonDesProgrammes.lean\#L57}

\begin{theoreme}
Après \texttt{n} tours à partir de \texttt{(0, 0)}, le compteur vaut \texttt{n} et l'accumulateur vaut
\texttt{n(n+1)/2} : le programme calcule bien la somme des \texttt{n} premiers entiers.

L'invariant est l'égalité \texttt{2s = i(i+1)}, qui évite la division ; la forme
\texttt{n(n+1)/2} s'en déduit à la fin.
\end{theoreme}
\begin{proof}
On démontre d'abord, par récurrence, un invariant plus fort que l'énoncé et sans
division : après \(m\) tours, le compteur vaut \(m\) et l'accumulateur \(s\) vérifie
\(2s = m(m+1)\).

Au tour \(0\), compteur et accumulateur sont nuls. Si l'invariant vaut après \(k\)
tours, un tour de plus porte le compteur à \(k+1\) et l'accumulateur à \(s + k + 1\), et
\[ 2(s + k + 1) = 2s + 2k + 2 = k(k+1) + 2k + 2 = (k+1)(k+2) . \]

L'énoncé s'en déduit : le produit \(n(n+1)\) étant pair, l'accumulateur en est la
moitié, c'est-à-dire \(n(n+1)/2\).

Écrire l'invariant sans division n'est pas une coquetterie. Sur les entiers, la division
est tronquée ; une récurrence menée directement sur \(n(n+1)/2\) obligerait à justifier
à chaque pas que la troncature ne fait rien.
\end{proof}
\source{https://github.com/Commutator-IO/learn-lean/blob/main/courses/02-lycee/11-informatique-nsi/CorrectionEtTerminaisonDesProgrammes.lean\#L64}{CorrectionEtTerminaisonDesProgrammes.lean\#L64}

\subsection{Division euclidienne par soustractions successives}

\begin{definition}
La division par soustractions successives : tant que le dividende est au moins
égal au diviseur, on retranche le diviseur et l'on compte un quotient de plus.

La récursion termine parce que \texttt{a - b < a} dès que \texttt{0 < b \(\le\) a} : la valeur du
dividende est le variant de la boucle.
\end{definition}
\source{https://github.com/Commutator-IO/learn-lean/blob/main/courses/02-lycee/11-informatique-nsi/CorrectionEtTerminaisonDesProgrammes.lean\#L98}{CorrectionEtTerminaisonDesProgrammes.lean\#L98}

\begin{theoreme}
Le programme rend bien le quotient et le reste : le quotient multiplié par le
diviseur, plus le reste, redonne le dividende, et le reste est strictement plus petit
que le diviseur.
\end{theoreme}
\begin{proof}
Par récurrence forte sur le dividende \(a\).

Si \(b > a\), la boucle ne s'exécute pas : le quotient est \(0\) et le reste \(a\), qui
est bien inférieur à \(b\), et \(0 \times b + a = a\).

Sinon \(b \le a\), et l'appel se fait sur \(a - b\), strictement plus petit que \(a\)
puisque \(b > 0\) : c'est le variant de la boucle. Par hypothèse de récurrence il rend
un quotient \(q\) et un reste \(r < b\) avec \(qb + r = a - b\). Le programme rend alors
\(q + 1\) et \(r\), et
\[ (q+1)b + r = (qb + r) + b = (a - b) + b = a . \]

L'hypothèse \(b > 0\) est indispensable : sans elle la boucle ne terminerait pas, et
l'énoncé n'aurait pas de sens.
\end{proof}
\source{https://github.com/Commutator-IO/learn-lean/blob/main/courses/02-lycee/11-informatique-nsi/CorrectionEtTerminaisonDesProgrammes.lean\#L109}{CorrectionEtTerminaisonDesProgrammes.lean\#L109}

\subsection{Deux programmes pour une même fonction}

\begin{definition}
La même somme, écrite récursivement : la structure du programme n'a rien à voir
avec la précédente \textemdash{} pas de boucle, pas d'accumulateur.
\end{definition}
\source{https://github.com/Commutator-IO/learn-lean/blob/main/courses/02-lycee/11-informatique-nsi/CorrectionEtTerminaisonDesProgrammes.lean\#L134}{CorrectionEtTerminaisonDesProgrammes.lean\#L134}

\begin{theoreme}
Deux programmes de structures différentes calculent la même fonction : la boucle
avec accumulateur et la définition récursive rendent la même valeur pour toute
entrée. C'est le sens qu'on donne à \og{} ces deux programmes sont corrects \fg{} : ils ne
se ressemblent pas, ils calculent la même chose.
\end{theoreme}
\begin{proof}
Par récurrence sur \(n\). En \(0\), les deux programmes rendent \(0\).

Au rang \(k+1\), la version récursive ajoute \(k+1\) à ce qu'elle a calculé au rang
\(k\), tandis que la boucle ajoute à son accumulateur la valeur de son compteur
augmentée de \(1\). Or le compteur vaut \(k\) après \(k\) tours, comme l'établit
l'invariant du programme précédent : les deux ajouts coïncident, et l'hypothèse de
récurrence donne l'égalité des accumulateurs.

C'est ce qu'on entend en disant que deux programmes sont corrects : ils ne se
ressemblent pas — l'un boucle, l'autre s'appelle lui-même — et pourtant ils calculent
la même fonction.
\end{proof}
\source{https://github.com/Commutator-IO/learn-lean/blob/main/courses/02-lycee/11-informatique-nsi/CorrectionEtTerminaisonDesProgrammes.lean\#L142}{CorrectionEtTerminaisonDesProgrammes.lean\#L142}

\subsection{Indécidabilité du problème de l'arrêt}

\begin{theoreme}
Le problème de l'arrêt est indécidable.

On se donne un langage de programmation abstrait : \texttt{arrete p e} signifie que le
programme \texttt{p} s'arrête sur l'entrée \texttt{e}, et tout programme peut être donné en entrée
à un autre \textemdash{} c'est le cas dans un vrai langage, où un programme est un fichier
texte.

L'hypothèse \texttt{hdiagonal} est l'unique chose qu'on demande au langage : supposer qu'un
programme décide l'arrêt permet d'en écrire un autre, le programme diagonal, qui
s'arrête sur \texttt{p} exactement quand \texttt{p} ne s'arrête pas sur lui-même \textemdash{} il suffit
d'appeler le décideur et de boucler exprès quand il répond \og{} oui \fg{}. La contradiction
est alors immédiate, en appliquant le programme diagonal à lui-même.

Ce que cet énoncé ne formalise pas : que ce programme diagonal soit effectivement
écrivable. C'est la partie qui demande un modèle de calcul \textemdash{} machines de Turing,
fonctions récursives \textemdash{} et qui dépasse le programme de terminale ; elle est ici
portée par l'hypothèse, énoncée en toutes lettres plutôt que passée sous silence.
\end{theoreme}
\begin{proof}
Il suffit d'appliquer le programme diagonal à lui-même. Par hypothèse, il s'arrête sur
un programme \(p\) exactement quand \(p\) ne s'arrête pas sur lui-même ; pour \(p\) égal
au programme diagonal, cela donne
\[ \text{le diagonal s'arrête sur lui-même} \iff
   \text{le diagonal ne s'arrête pas sur lui-même} , \]
ce qui est contradictoire.

Toute la difficulté du théorème est dans l'hypothèse, non dans cette contradiction. Ce
qu'il faudrait établir, et qui n'est pas fait ici, c'est qu'un programme décidant
l'arrêt permettrait d'écrire effectivement ce programme diagonal — il suffirait de
l'appeler et de boucler exprès quand il répond « oui ». Le rendre rigoureux demande un
modèle de calcul, machines de Turing ou fonctions récursives, que le programme de
terminale n'introduit pas. L'hypothèse le dit en toutes lettres plutôt que de le passer
sous silence.
\end{proof}
\source{https://github.com/Commutator-IO/learn-lean/blob/main/courses/02-lycee/11-informatique-nsi/CorrectionEtTerminaisonDesProgrammes.lean\#L174}{CorrectionEtTerminaisonDesProgrammes.lean\#L174}

\section{RecursiviteEtDiviserPourRegner}

\noindent
Lycée \textemdash{} spécialité NSI, section \og{} Récursivité et diviser pour régner \fg{}.
Une fonction récursive s'appelle elle-même : rien ne garantit a priori qu'elle
s'arrête, ni qu'elle calcule ce qu'on croit. Ce fichier démontre les deux pour la
factorielle, Fibonacci et l'exponentiation rapide, puis compare les coûts \textemdash{} la
différence entre le Fibonacci récursif naïf et sa version itérative est
exponentielle, et couper un problème en deux moitiés donne le coût \texttt{n log₂ n}.
Énoncés et démonstrations en français : voir Informatique.tex.

\subsection{Factorielle récursive}

\begin{definition}
La factorielle, écrite récursivement. La récursion porte sur le constructeur de
l'entier : chaque appel se fait sur \texttt{n}, prédécesseur de \texttt{n + 1}. C'est ce qui fait
que la fonction termine, et Lean l'accepte sans qu'on ait à le justifier.
\end{definition}
\source{https://github.com/Commutator-IO/learn-lean/blob/main/courses/02-lycee/11-informatique-nsi/RecursiviteEtDiviserPourRegner.lean\#L19}{RecursiviteEtDiviserPourRegner.lean\#L19}

\begin{theoreme}
La fonction récursive calcule bien la factorielle : son résultat est le produit
des entiers de \texttt{1} à \texttt{n}. Elle ne rend jamais \texttt{0}, ce qui autorise à diviser par
elle.
\end{theoreme}
\begin{proof}
Par récurrence. En \(0\), la fonction rend \(1\), qui est le produit vide.

Au rang \(k+1\), elle rend \((k+1) \times \mathrm{factorielle}(k)\) ; l'hypothèse de
récurrence identifie le second facteur au produit des entiers de \(1\) à \(k\), et
ajouter le facteur \(k+1\) donne le produit jusqu'à \(k+1\). La positivité suit de même,
un produit de facteurs non nuls étant non nul.

La terminaison ne demande rien : l'appel se fait sur \(k\), prédécesseur immédiat de
\(k+1\). C'est le cas le plus simple, et Lean l'accepte sans qu'on ait à écrire quoi que
ce soit.
\end{proof}
\source{https://github.com/Commutator-IO/learn-lean/blob/main/courses/02-lycee/11-informatique-nsi/RecursiviteEtDiviserPourRegner.lean\#L26}{RecursiviteEtDiviserPourRegner.lean\#L26}

\subsection{Suite de Fibonacci}

\begin{definition}
Fibonacci, écrit récursivement comme la définition mathématique le dit : chaque
terme est la somme des deux précédents. Deux appels par terme.
\end{definition}
\source{https://github.com/Commutator-IO/learn-lean/blob/main/courses/02-lycee/11-informatique-nsi/RecursiviteEtDiviserPourRegner.lean\#L40}{RecursiviteEtDiviserPourRegner.lean\#L40}

\begin{definition}
Fibonacci, écrit itérativement : on garde deux termes consécutifs et l'on avance
d'un cran à chaque tour. Un seul parcours.
\end{definition}
\source{https://github.com/Commutator-IO/learn-lean/blob/main/courses/02-lycee/11-informatique-nsi/RecursiviteEtDiviserPourRegner.lean\#L47}{RecursiviteEtDiviserPourRegner.lean\#L47}

\begin{theoreme}
Les deux versions calculent la même valeur.

La récurrence porte sur l'invariant de la boucle, plus fort que l'énoncé : le couple
gardé en mémoire est exactement \texttt{(F(n), F(n+1))}. C'est un cas typique où l'on ne
peut pas démontrer l'énoncé tel quel, et où il faut d'abord dire ce que le programme
conserve.
\end{theoreme}
\begin{proof}
Par récurrence, mais pas sur l'énoncé : on démontre l'invariant de la boucle, qui est
plus fort — le couple gardé en mémoire après \(n\) tours est exactement
\(\bigl(F(n),\, F(n+1)\bigr)\).

Au tour \(0\), ce couple est \((0, 1)\), qui est bien \(\bigl(F(0), F(1)\bigr)\). Si
l'invariant vaut au tour \(k\), un tour de plus rend
\(\bigl(F(k+1),\, F(k) + F(k+1)\bigr)\), c'est-à-dire \(\bigl(F(k+1), F(k+2)\bigr)\) par
définition de la suite.

L'énoncé s'obtient en ne gardant que la première composante. C'est un cas typique où
l'énoncé qu'on veut ne se démontre pas tel quel : il faut d'abord dire ce que le
programme conserve.
\end{proof}
\source{https://github.com/Commutator-IO/learn-lean/blob/main/courses/02-lycee/11-informatique-nsi/RecursiviteEtDiviserPourRegner.lean\#L57}{RecursiviteEtDiviserPourRegner.lean\#L57}

\subsection{Coût des deux versions}

\begin{definition}
Le nombre d'appels que fait le Fibonacci récursif naïf : un pour lui-même, plus
ceux de ses deux appels.
\end{definition}
\source{https://github.com/Commutator-IO/learn-lean/blob/main/courses/02-lycee/11-informatique-nsi/RecursiviteEtDiviserPourRegner.lean\#L69}{RecursiviteEtDiviserPourRegner.lean\#L69}

\begin{definition}
Le nombre de tours de la version itérative : un par cran.
\end{definition}
\source{https://github.com/Commutator-IO/learn-lean/blob/main/courses/02-lycee/11-informatique-nsi/RecursiviteEtDiviserPourRegner.lean\#L75}{RecursiviteEtDiviserPourRegner.lean\#L75}

\begin{theoreme}
Le nombre d'appels croît avec \texttt{n} : chaque appel supplémentaire ne peut
qu'ajouter du travail. C'est l'étape qu'il faut avoir démontrée avant de minorer le
coût.
\end{theoreme}
\begin{proof}
Par récurrence forte, en distinguant les trois formes que peut prendre \(n\).

Pour \(n = 0\) et \(n = 1\), le nombre d'appels vaut \(1\), contre \(1\) et \(3\)
respectivement.

Pour \(n = k+2\), le nombre d'appels vaut \(1 + a(k) + a(k+1)\), et celui de \(k+3\)
vaut \(1 + a(k+1) + a(k+2)\) ; l'hypothèse de récurrence appliquée à \(k\) donne
\(a(k) \le a(k+1)\), ce qui conclut.

Cette croissance est ce qu'il faut avoir en main avant de minorer le coût.
\end{proof}
\source{https://github.com/Commutator-IO/learn-lean/blob/main/courses/02-lycee/11-informatique-nsi/RecursiviteEtDiviserPourRegner.lean\#L82}{RecursiviteEtDiviserPourRegner.lean\#L82}

\begin{theoreme}
Le coût du Fibonacci récursif naïf est exponentiel, celui de la version
itérative est linéaire.

La minoration \texttt{2\textasciicircum{}(n/2) \(\le\) appels(n)} suffit à conclure : le nombre d'appels double au
moins tous les deux crans, alors que la boucle n'en fait que \texttt{n}. Pour \texttt{n = 40}, cela
sépare le million de la quarantaine.
\end{theoreme}
\begin{proof}
Deux démonstrations séparées.

Pour la version récursive, récurrence forte sur \(n\). Pour \(n = 0\) et \(n = 1\), le
nombre d'appels vaut \(1\), et \(2^0 = 1\). Pour \(n = k+2\),
\[ a(k+2) = 1 + a(k) + a(k+1) \ge 2\,a(k)
   \ge 2 \times 2^{\lfloor k/2 \rfloor} = 2^{\lfloor (k+2)/2 \rfloor} , \]
en utilisant d'abord la croissance établie plus haut, puis l'hypothèse de récurrence.

Pour la version itérative, récurrence immédiate : un tour par cran.

L'écart est celui du million à la quarantaine pour \(n = 40\). Il ne vient pas d'une
différence d'idée mais d'une différence d'écriture : la version naïve recalcule
indéfiniment les mêmes termes.
\end{proof}
\source{https://github.com/Commutator-IO/learn-lean/blob/main/courses/02-lycee/11-informatique-nsi/RecursiviteEtDiviserPourRegner.lean\#L100}{RecursiviteEtDiviserPourRegner.lean\#L100}

\subsection{Exponentiation rapide}

\begin{definition}
L'exponentiation rapide : au lieu de multiplier \texttt{n} fois par \texttt{a}, on élève au
carré la puissance de l'exposant moitié. La récursion se fait sur \texttt{n / 2}, qui est
strictement plus petit que \texttt{n} dès que \texttt{n} n'est pas nul : c'est le variant qui
justifie la terminaison, et Lean le réclame explicitement.
\end{definition}
\source{https://github.com/Commutator-IO/learn-lean/blob/main/courses/02-lycee/11-informatique-nsi/RecursiviteEtDiviserPourRegner.lean\#L125}{RecursiviteEtDiviserPourRegner.lean\#L125}

\begin{definition}
Le nombre d'élévations au carré que fait l'exponentiation rapide : une par
appel.
\end{definition}
\source{https://github.com/Commutator-IO/learn-lean/blob/main/courses/02-lycee/11-informatique-nsi/RecursiviteEtDiviserPourRegner.lean\#L134}{RecursiviteEtDiviserPourRegner.lean\#L134}

\begin{theoreme}
L'exponentiation rapide calcule bien \texttt{aⁿ}.

Tout repose sur l'écriture \texttt{n = 2⌊n/2⌋ + (n mod 2)}, donc
\texttt{aⁿ = (a\textasciicircum{}⌊n/2⌋)² \(\times\) a\textasciicircum{}(n mod 2)} : c'est le carré qui fait tout le travail, et le
facteur \texttt{a} supplémentaire ne sert que pour un exposant impair.
\end{theoreme}
\begin{proof}
Par récurrence forte sur l'exposant. Si \(n = 0\), la fonction rend \(1\).

Sinon elle calcule \(r = a^{\lfloor n/2 \rfloor}\), par hypothèse de récurrence :
l'appel se fait sur \(\lfloor n/2 \rfloor\), strictement plus petit que \(n\). Tout
tient alors à l'écriture
\[ n = 2\left\lfloor n/2 \right\rfloor + (n \bmod 2) . \]
Si \(n\) est pair, \(a^n = r^2\) ; s'il est impair, \(a^n = a \times r^2\). C'est
exactement ce que rend la fonction dans chacun des deux cas.
\end{proof}
\source{https://github.com/Commutator-IO/learn-lean/blob/main/courses/02-lycee/11-informatique-nsi/RecursiviteEtDiviserPourRegner.lean\#L144}{RecursiviteEtDiviserPourRegner.lean\#L144}

\begin{theoreme}
L'exponentiation rapide fait au plus \texttt{⌊log₂ n⌋ + 1} élévations au carré : chaque
appel divise l'exposant par deux, et l'on ne peut diviser par deux qu'un nombre
logarithmique de fois avant d'atteindre zéro.
\end{theoreme}
\begin{proof}
Par récurrence forte. Pour \(n = 1\), un seul appel, et \(\log_2 1 + 1 = 1\).

Pour \(n \ge 2\), l'exposant moitié est encore non nul, et
\(\log_2 \lfloor n/2 \rfloor = \log_2 n - 1\). L'hypothèse de récurrence majore le coût
de l'appel par \(\log_2 n - 1 + 1\), et il reste à compter l'élévation au carré du tour
courant, d'où \(\log_2 n + 1\).

Chaque appel divise l'exposant par deux, et l'on ne peut diviser par deux qu'un nombre
logarithmique de fois avant d'atteindre zéro : calculer \(a^{1000}\) demande dix
élévations au carré, non mille multiplications.
\end{proof}
\source{https://github.com/Commutator-IO/learn-lean/blob/main/courses/02-lycee/11-informatique-nsi/RecursiviteEtDiviserPourRegner.lean\#L163}{RecursiviteEtDiviserPourRegner.lean\#L163}

\subsection{Tours de Hanoï}

\begin{definition}
Le nombre de déplacements de la solution récursive : pour déplacer \texttt{n + 1}
disques, on déplace les \texttt{n} plus petits sur le piquet libre, on déplace le grand, on
remet les \texttt{n} petits dessus.
\end{definition}
\source{https://github.com/Commutator-IO/learn-lean/blob/main/courses/02-lycee/11-informatique-nsi/RecursiviteEtDiviserPourRegner.lean\#L184}{RecursiviteEtDiviserPourRegner.lean\#L184}

\begin{theoreme}
La solution récursive demande \texttt{2ⁿ − 1} déplacements. L'égalité est écrite sous la
forme \texttt{hanoi n + 1 = 2ⁿ} pour éviter la soustraction tronquée des entiers naturels,
où \texttt{0 − 1} vaudrait \texttt{0}.
\end{theoreme}
\begin{proof}
Par récurrence. En \(0\), aucun déplacement, et \(0 + 1 = 2^0\).

Au rang \(k+1\), la solution récursive déplace les \(k\) plus petits disques sur le
piquet libre, puis le grand, puis à nouveau les \(k\) petits : cela fait
\(2h(k) + 1\) déplacements. L'hypothèse de récurrence donne \(h(k) + 1 = 2^k\), d'où
\[ h(k+1) + 1 = 2h(k) + 2 = 2\bigl(h(k) + 1\bigr) = 2 \times 2^k = 2^{k+1} . \]

L'égalité est écrite sous la forme \(h(n) + 1 = 2^n\) pour éviter la soustraction
tronquée des entiers naturels, où \(0 - 1\) vaudrait \(0\).
\end{proof}
\source{https://github.com/Commutator-IO/learn-lean/blob/main/courses/02-lycee/11-informatique-nsi/RecursiviteEtDiviserPourRegner.lean\#L191}{RecursiviteEtDiviserPourRegner.lean\#L191}

\begin{figure}[h]
\centering
% Construction : trois piquets en x = -3, 0, 3, socle en y = 0, hauteur 2.2.
%   Trois disques empilés sur le premier, de demi-largeurs 1.1, 0.8, 0.5,
%   aux hauteurs 0.15, 0.55, 0.95 (épaisseur 0.32).
% SVG : x = 150 + 34x, y = 120 - 34y.
\begin{tikzpicture}[scale=1, >=stealth]
  \draw (-4.3, 0) -- (4.3, 0);
  \foreach \x in {-3, 0, 3} \draw (\x, 0) -- (\x, 2.2);
  \foreach \w/\y in {1.1/0.15, 0.8/0.55, 0.5/0.95}
    \draw (-3-\w, \y) rectangle (-3+\w, \y+0.32);
  \draw[->] (-1.6, 1.9) to[bend left=20] (-0.2, 1.9);
  \node[above, font=\footnotesize] at (-0.9, 2.1) {$2^n - 1$};
\end{tikzpicture}

\caption{Les tours de Hanoï. Déplacer la pile entière demande d'abord de dégager les disques du dessus, puis de les remettre : d'où le doublement à chaque disque, et les $2^n - 1$ déplacements.}
\end{figure}

\begin{theoreme}
Aucune solution ne fait mieux.

L'hypothèse \texttt{hgrand} est la seule observation qu'on demande sur les déplacements :
pour déplacer \texttt{n + 1} disques, il faut bien à un moment déplacer le plus grand, et à
cet instant les \texttt{n} autres sont tous sur le piquet restant \textemdash{} il a donc fallu les y
amener, puis les ramener ensuite, d'où au moins \texttt{2f(n) + 1} déplacements. La
récurrence en tire la minoration pour tout \texttt{n}.

Ce que cet énoncé ne formalise pas : cette observation elle-même, qui demanderait de
modéliser les configurations et les déplacements licites. Elle est portée par
l'hypothèse, écrite en toutes lettres.
\end{theoreme}
\begin{proof}
Par récurrence. En \(0\), la solution récursive ne fait aucun déplacement.

Au rang \(k+1\), l'hypothèse de récurrence donne \(h(k) \le f(k)\), et l'hypothèse sur
\(f\) donne \(2f(k) + 1 \le f(k+1)\), d'où
\[ h(k+1) = 2h(k) + 1 \le 2f(k) + 1 \le f(k+1) . \]

L'hypothèse \(2f(n) + 1 \le f(n+1)\) est l'observation qui fait tout le travail : pour
déplacer \(n+1\) disques il faut bien, à un moment, déplacer le plus grand, et à cet
instant les \(n\) autres sont tous sur le piquet restant — il a donc fallu les y
amener, puis les ramener.

Ce que cette démonstration ne formalise pas, c'est cette observation elle-même, qui
demanderait de modéliser les configurations et les déplacements licites. Elle est portée
par l'hypothèse, écrite en toutes lettres.
\end{proof}
\source{https://github.com/Commutator-IO/learn-lean/blob/main/courses/02-lycee/11-informatique-nsi/RecursiviteEtDiviserPourRegner.lean\#L209}{RecursiviteEtDiviserPourRegner.lean\#L209}

\subsection{Diviser pour régner}

\begin{definition}
Le coût d'un algorithme qui coupe le problème en deux moitiés, les traite, puis
recolle en temps linéaire \textemdash{} c'est le schéma du tri fusion. La variable est ici le
nombre d'étages \texttt{k}, c'est-à-dire la taille \texttt{n = 2ᵏ}.
\end{definition}
\source{https://github.com/Commutator-IO/learn-lean/blob/main/courses/02-lycee/11-informatique-nsi/RecursiviteEtDiviserPourRegner.lean\#L223}{RecursiviteEtDiviserPourRegner.lean\#L223}

\begin{theoreme}
Le coût vaut \texttt{n log₂ n} : pour \texttt{n = 2ᵏ}, il vaut \texttt{k \(\cdot\) 2ᵏ}.

C'est l'énoncé \texttt{T(n) = 2T(n/2) + n} résolu. On le lit ainsi : il y a \texttt{k = log₂ n}
étages de découpe, et chaque étage coûte \texttt{n} au total, puisque les recollements d'un
même étage portent ensemble sur les \texttt{n} éléments.
\end{theoreme}
\begin{proof}
Par récurrence sur le nombre d'étages. À zéro étage, il n'y a rien à faire.

À \(m+1\) étages, le coût vaut \(2T(m) + 2^{m+1}\), et l'hypothèse de récurrence donne
\(T(m) = m\,2^m\), d'où
\[ T(m+1) = 2m\,2^m + 2^{m+1} = m\,2^{m+1} + 2^{m+1} = (m+1)\,2^{m+1} . \]

C'est la récurrence \(T(n) = 2T(n/2) + n\) résolue, pour \(n = 2^k\) : il y a
\(k = \log_2 n\) étages de découpe, et chaque étage coûte \(n\) au total, puisque les
recollements d'un même étage portent ensemble sur les \(n\) éléments.
\end{proof}
\source{https://github.com/Commutator-IO/learn-lean/blob/main/courses/02-lycee/11-informatique-nsi/RecursiviteEtDiviserPourRegner.lean\#L232}{RecursiviteEtDiviserPourRegner.lean\#L232}

\section{StructuresDeDonnees}

\noindent
Lycée \textemdash{} spécialité NSI, section \og{} Structures de données \fg{}.
Une structure de données n'est pas un rangement : c'est un contrat. La pile promet que
le dernier entré sort le premier, la file que c'est le premier, l'arbre binaire de
recherche que le parcours infixe rend les clés triées. Ce fichier écrit ces structures
et démontre leurs contrats \textemdash{} y compris celui, moins évident, de la file simulée par
deux piles.
Énoncés et démonstrations en français : voir Informatique.tex.

\subsection{Piles}

\begin{definition}
Empiler : le nouvel élément est posé au sommet.
\end{definition}
\source{https://github.com/Commutator-IO/learn-lean/blob/main/courses/02-lycee/11-informatique-nsi/StructuresDeDonnees.lean\#L19}{StructuresDeDonnees.lean\#L19}

\begin{definition}
Dépiler : on rend le sommet et la pile privée de ce sommet, ou rien si la pile est
vide.
\end{definition}
\source{https://github.com/Commutator-IO/learn-lean/blob/main/courses/02-lycee/11-informatique-nsi/StructuresDeDonnees.lean\#L23}{StructuresDeDonnees.lean\#L23}

\begin{theoreme}
Le contrat de la pile : dernier entré, premier sorti. Dépiler juste après avoir
empilé rend l'élément qu'on vient de poser et l'état d'avant ; et si l'on empile deux
éléments, c'est le second qui ressort d'abord.
\end{theoreme}
\begin{proof}
Les deux égalités sont vraies par définition : empiler pose l'élément en tête, dépiler
la retire. Il n'y a rien à calculer — et c'est bien le propos : le contrat de la pile
est exactement sa définition.
\end{proof}
\source{https://github.com/Commutator-IO/learn-lean/blob/main/courses/02-lycee/11-informatique-nsi/StructuresDeDonnees.lean\#L30}{StructuresDeDonnees.lean\#L30}

\subsection{Files}

\begin{definition}
Enfiler : le nouvel élément est mis à la fin.
\end{definition}
\source{https://github.com/Commutator-IO/learn-lean/blob/main/courses/02-lycee/11-informatique-nsi/StructuresDeDonnees.lean\#L38}{StructuresDeDonnees.lean\#L38}

\begin{definition}
Défiler : on rend la tête, c'est-à-dire l'élément qui attend depuis le plus
longtemps.
\end{definition}
\source{https://github.com/Commutator-IO/learn-lean/blob/main/courses/02-lycee/11-informatique-nsi/StructuresDeDonnees.lean\#L42}{StructuresDeDonnees.lean\#L42}

\begin{theoreme}
Le contrat de la file : premier entré, premier sorti. Enfiler dans une file vide
puis défiler rend l'élément ; enfiler dans une file non vide ne change pas qui sort en
premier \textemdash{} le nouveau venu prend la queue.
\end{theoreme}
\begin{proof}
La première égalité est vraie par définition : enfiler dans la file vide donne \([x]\),
dont la tête est \(x\).

La seconde demande un calcul. Enfiler \(x\) dans \(y :: f\) donne la concaténation de
\(y :: f\) et de \([x]\), c'est-à-dire \(y\) suivi de la concaténation de \(f\) et de
\([x]\). Défiler rend donc \(y\) — le nouveau venu n'a pas doublé la file — et laisse
\(f\) augmentée de \(x\) en queue.
\end{proof}
\source{https://github.com/Commutator-IO/learn-lean/blob/main/courses/02-lycee/11-informatique-nsi/StructuresDeDonnees.lean\#L49}{StructuresDeDonnees.lean\#L49}

\subsection{File par deux piles}

\begin{definition}
Une file simulée par deux piles : on empile les arrivées sur la pile d'entrée, et
l'on sert depuis la pile de sortie ; quand celle-ci se vide, on y bascule l'entrée
retournée. C'est la construction classique, et elle n'a rien d'évident : la file
apparaît alors que ni l'une ni l'autre des deux piles ne se comporte comme une file.
\end{definition}
\source{https://github.com/Commutator-IO/learn-lean/blob/main/courses/02-lycee/11-informatique-nsi/StructuresDeDonnees.lean\#L62}{StructuresDeDonnees.lean\#L62}

\begin{definition}
Enfiler dans la structure : on empile sur la pile d'entrée.
\end{definition}
\source{https://github.com/Commutator-IO/learn-lean/blob/main/courses/02-lycee/11-informatique-nsi/StructuresDeDonnees.lean\#L67}{StructuresDeDonnees.lean\#L67}

\begin{definition}
Défiler dans la structure : on sert la pile de sortie ; si elle est vide, on y
bascule d'abord la pile d'entrée retournée.
\end{definition}
\source{https://github.com/Commutator-IO/learn-lean/blob/main/courses/02-lycee/11-informatique-nsi/StructuresDeDonnees.lean\#L72}{StructuresDeDonnees.lean\#L72}

\begin{definition}
Le contenu de la structure, vu comme une file ordinaire : ce qui attend dans la
pile de sortie, puis ce qui attend dans la pile d'entrée, retourné. C'est la fonction
d'abstraction \textemdash{} c'est elle qui donne un sens à \og{} se comporter comme une file \fg{}.
\end{definition}
\source{https://github.com/Commutator-IO/learn-lean/blob/main/courses/02-lycee/11-informatique-nsi/StructuresDeDonnees.lean\#L83}{StructuresDeDonnees.lean\#L83}

\begin{theoreme}
La file par deux piles se comporte bien comme une file : ses deux opérations
correspondent, à travers la fonction d'abstraction, à celles de la file ordinaire.
C'est ce qu'il faut démontrer pour avoir le droit de remplacer l'une par l'autre.
\end{theoreme}
\begin{proof}
On distingue les cas selon l'état des deux piles. Notons \(\overline{\ell}\) le
retournement de \(\ell\).

Pour l'enfilement, le contenu de la structure obtenue est
\[ \mathrm{sortie} \mathbin{+\!\!+} \overline{x :: \mathrm{entree}}
   = \mathrm{sortie} \mathbin{+\!\!+}
     \bigl(\overline{\mathrm{entree}} \mathbin{+\!\!+} [x]\bigr)
   = \bigl(\mathrm{sortie} \mathbin{+\!\!+} \overline{\mathrm{entree}}\bigr)
     \mathbin{+\!\!+} [x] , \]
c'est-à-dire le contenu d'avant, augmenté de \(x\) en queue.

Pour le défilement, deux cas. Si la pile de sortie n'est pas vide, sa tête est aussi la
tête du contenu, et les deux opérations la retirent. Si elle est vide, le contenu est
l'entrée retournée : ou bien celle-ci est vide aussi, et les deux côtés rendent
« rien » ; ou bien elle s'écrit \(z :: t\), et la structure bascule \(t\) dans la pile de
sortie, dont le contenu est alors \(t\) — exactement ce que rend le défilement de la
file ordinaire.

C'est ce basculement qui rend la construction intéressante : ni l'une ni l'autre des
deux piles ne se comporte comme une file, et pourtant leur composition, oui.
\end{proof}
\source{https://github.com/Commutator-IO/learn-lean/blob/main/courses/02-lycee/11-informatique-nsi/StructuresDeDonnees.lean\#L88}{StructuresDeDonnees.lean\#L88}

\begin{figure}[h]
\centering
% Construction : deux piles côte à côte, entrée à gauche (x = -2.4), sortie à
%   droite (x = 1.6), largeur 1.6, socle en y = 0, éléments d'épaisseur 0.55.
%   Entrée empilée d, c, b de bas en haut ; sortie vide ; flèche de bascule.
% SVG : x = 150 + 34x, y = 122 - 34y.
\begin{tikzpicture}[scale=1, >=stealth, every node/.style={font=\small}]
  \draw (-3.2, 0) -- (-1.6, 0) -- (-1.6, 2.2); \draw (-3.2, 0) -- (-3.2, 2.2);
  \draw (0.8, 0) -- (2.4, 0) -- (2.4, 2.2); \draw (0.8, 0) -- (0.8, 2.2);
  \foreach \e/\y in {d/0.05, c/0.65, b/1.25}
    \draw (-3.1, \y) rectangle node {$\e$} (-1.7, \y+0.55);
  \draw[->] (-1.4, 1.9) to[bend left=25] (0.6, 1.9);
  \node[above, font=\footnotesize] at (-0.4, 2.15) {bascule, retournée};
  \node[below] at (-2.4, -0.15) {entrée};
  \node[below] at (1.6, -0.15) {sortie};
\end{tikzpicture}

\caption{La file simulée par deux piles. Les arrivées s'empilent à l'entrée ; quand la pile de sortie se vide, on y bascule l'entrée, ce qui la retourne \textemdash{} et rend au premier arrivé sa place de premier servi.}
\end{figure}

\subsection{Listes chaînées}

\begin{definition}
La longueur d'une liste chaînée, comptée maillon par maillon.
\end{definition}
\source{https://github.com/Commutator-IO/learn-lean/blob/main/courses/02-lycee/11-informatique-nsi/StructuresDeDonnees.lean\#L104}{StructuresDeDonnees.lean\#L104}

\begin{definition}
La concaténation de deux listes chaînées : on recopie la première devant la
seconde.
\end{definition}
\source{https://github.com/Commutator-IO/learn-lean/blob/main/courses/02-lycee/11-informatique-nsi/StructuresDeDonnees.lean\#L110}{StructuresDeDonnees.lean\#L110}

\begin{definition}
Le parcours d'une liste chaînée : on visite les maillons du premier au dernier.
\end{definition}
\source{https://github.com/Commutator-IO/learn-lean/blob/main/courses/02-lycee/11-informatique-nsi/StructuresDeDonnees.lean\#L115}{StructuresDeDonnees.lean\#L115}

\begin{theoreme}
Les deux propriétés qu'on attend d'une liste chaînée : la longueur d'une
concaténation est la somme des longueurs, et le parcours voit exactement les éléments
de la liste \textemdash{} il n'en oublie aucun et n'en invente aucun.
\end{theoreme}
\begin{proof}
Deux récurrences sur la première liste.

Pour la longueur : concaténer la liste vide devant \(l_2\) rend \(l_2\), et la somme des
longueurs vaut celle de \(l_2\). Sinon la concaténation ajoute un maillon devant, et les
deux membres augmentent de \(1\).

Pour le parcours : la liste vide se parcourt en la liste vide ; sinon le parcours visite
la tête puis parcourt la queue, et l'hypothèse de récurrence donne le reste. Le parcours
n'oublie aucun maillon et n'en invente aucun.
\end{proof}
\source{https://github.com/Commutator-IO/learn-lean/blob/main/courses/02-lycee/11-informatique-nsi/StructuresDeDonnees.lean\#L122}{StructuresDeDonnees.lean\#L122}

\subsection{Arbres binaires}

\begin{definition}
Un arbre binaire : soit vide, soit un nœud portant une valeur et deux
sous-arbres.
\end{definition}
\source{https://github.com/Commutator-IO/learn-lean/blob/main/courses/02-lycee/11-informatique-nsi/StructuresDeDonnees.lean\#L138}{StructuresDeDonnees.lean\#L138}

\begin{definition}
Le nombre de nœuds.
\end{definition}
\source{https://github.com/Commutator-IO/learn-lean/blob/main/courses/02-lycee/11-informatique-nsi/StructuresDeDonnees.lean\#L143}{StructuresDeDonnees.lean\#L143}

\begin{definition}
La hauteur, comptée en nombre de niveaux : l'arbre vide a la hauteur \texttt{0}, un arbre
réduit à un nœud la hauteur \texttt{1}. C'est la convention des tableaux du programme ; avec
celle qui compte les arêtes, tous les énoncés se décalent de un.
\end{definition}
\source{https://github.com/Commutator-IO/learn-lean/blob/main/courses/02-lycee/11-informatique-nsi/StructuresDeDonnees.lean\#L150}{StructuresDeDonnees.lean\#L150}

\begin{theoreme}
Un arbre de hauteur \texttt{h} a au plus \texttt{2ʰ − 1} nœuds : chaque niveau ne peut au mieux
que doubler le nombre de nœuds du précédent. L'inégalité est écrite \texttt{n + 1 \(\le\) 2ʰ} pour
éviter la soustraction tronquée des entiers naturels.
\end{theoreme}
\begin{proof}
Par récurrence sur l'arbre. L'arbre vide a zéro nœud et la hauteur \(0\), et
\(0 + 1 \le 2^0\).

Un nœud de sous-arbres \(g\) et \(d\) a \(t(g) + 1 + t(d)\) nœuds et la hauteur
\(1 + \max(h(g), h(d))\). Par hypothèse de récurrence, \(t(g) + 1 \le 2^{h(g)}\) et
\(t(d) + 1 \le 2^{h(d)}\), tous deux majorés par \(2^{\max(h(g),\, h(d))}\). D'où
\[ \bigl(t(g) + 1 + t(d)\bigr) + 1 \le 2 \times 2^{\max(h(g),\, h(d))}
   = 2^{1 + \max(h(g),\, h(d))} . \]

La hauteur est comptée en niveaux : l'arbre vide a la hauteur \(0\), un arbre réduit à
un nœud la hauteur \(1\). Avec la convention qui compte les arêtes, la majoration
s'écrit \(2^{h+1} - 1\).
\end{proof}
\source{https://github.com/Commutator-IO/learn-lean/blob/main/courses/02-lycee/11-informatique-nsi/StructuresDeDonnees.lean\#L157}{StructuresDeDonnees.lean\#L157}

\begin{figure}[h]
\centering
% Construction (repère mathématique) : arbre complet de hauteur 3.
%   Niveau 1 : (0, 2). Niveau 2 : (-1.6, 1), (1.6, 1).
%   Niveau 3 : (-2.4, 0), (-0.8, 0), (0.8, 0), (2.4, 0).
% SVG : x = 130 + 40x, y = 30 + 40(2 - y).
\begin{tikzpicture}[scale=1, every node/.style={font=\small}]
  \tikzstyle{n}=[draw, circle, inner sep=0pt, minimum size=5mm]
  \node[n] (a) at (0, 2) {};
  \node[n] (b) at (-1.6, 1) {}; \node[n] (c) at (1.6, 1) {};
  \node[n] (d) at (-2.4, 0) {}; \node[n] (e) at (-0.8, 0) {};
  \node[n] (f) at (0.8, 0) {};  \node[n] (g) at (2.4, 0) {};
  \draw (a) -- (b); \draw (a) -- (c);
  \draw (b) -- (d); \draw (b) -- (e); \draw (c) -- (f); \draw (c) -- (g);
  \node[left] at (-3.1, 2) {$2^0 = 1$};
  \node[left] at (-3.1, 1) {$2^1 = 2$};
  \node[left] at (-3.1, 0) {$2^2 = 4$};
  \node[right, font=\footnotesize] at (2.9, 1) {$1 + 2 + 4 = 2^3 - 1$};
\end{tikzpicture}

\caption{Un arbre complet de hauteur $3$ : chaque niveau double le nombre de nœuds du précédent, et le total $1 + 2 + 4$ vaut $2^3 - 1$. Aucun arbre de cette hauteur n'en a davantage.}
\end{figure}

\begin{theoreme}
Conséquence directe : la hauteur d'un arbre à \texttt{n} nœuds vaut au moins
\texttt{log₂(n + 1)}. C'est la raison pour laquelle on cherche à équilibrer les arbres \textemdash{} on
ne peut pas descendre en dessous de cette hauteur, mais on peut s'en approcher.
\end{theoreme}
\begin{proof}
Il suffit d'appliquer le logarithme à la majoration précédente : de \(n + 1 \le 2^h\) on
tire \(\log_2(n+1) \le \log_2(2^h) = h\) par croissance.

C'est la raison pour laquelle on cherche à équilibrer les arbres. On ne peut pas
descendre en dessous de cette hauteur ; un arbre équilibré s'en approche, tandis qu'un
arbre dégénéré en peigne atteint l'autre extrême, avec une hauteur égale au nombre de
nœuds.
\end{proof}
\source{https://github.com/Commutator-IO/learn-lean/blob/main/courses/02-lycee/11-informatique-nsi/StructuresDeDonnees.lean\#L171}{StructuresDeDonnees.lean\#L171}

\begin{definition}
Un sous-arbre compte pour une arête s'il n'est pas vide.
\end{definition}
\source{https://github.com/Commutator-IO/learn-lean/blob/main/courses/02-lycee/11-informatique-nsi/StructuresDeDonnees.lean\#L178}{StructuresDeDonnees.lean\#L178}

\begin{definition}
Le nombre d'arêtes : une par sous-arbre non vide accroché à un nœud.
\end{definition}
\source{https://github.com/Commutator-IO/learn-lean/blob/main/courses/02-lycee/11-informatique-nsi/StructuresDeDonnees.lean\#L183}{StructuresDeDonnees.lean\#L183}

\begin{theoreme}
Un arbre non vide à \texttt{n} nœuds a \texttt{n − 1} arêtes : chaque nœud sauf la racine est
relié à son père par exactement une arête, et il n'y en a pas d'autre. C'est la
version enracinée de l'énoncé \og{} un arbre est un graphe connexe sans cycle, et il a
\texttt{n − 1} arêtes \fg{} ; la caractérisation par la connexité et l'absence de cycle demande
une théorie des graphes que ce fichier ne construit pas.
\end{theoreme}
\begin{proof}
Par récurrence sur l'arbre, en distinguant quatre cas selon que chaque sous-arbre est
vide ou non.

Si les deux le sont, l'arbre est réduit à un nœud : aucune arête, un nœud. Si un seul
l'est, disons le gauche, l'arbre a une arête vers le sous-arbre droit, plus les arêtes
de celui-ci, et l'hypothèse de récurrence appliquée à ce sous-arbre — non vide —
conclut. Si aucun ne l'est, on ajoute les deux arêtes et les deux hypothèses de
récurrence.

On peut le lire ainsi : chaque nœud sauf la racine est relié à son père par exactement
une arête, et il n'y en a pas d'autre.

L'énoncé est démontré pour les arbres enracinés, qui sont ceux que manipulent les
programmes. La caractérisation « un arbre est un graphe connexe sans cycle » demande une
théorie des graphes que ce chapitre ne construit pas.
\end{proof}
\source{https://github.com/Commutator-IO/learn-lean/blob/main/courses/02-lycee/11-informatique-nsi/StructuresDeDonnees.lean\#L192}{StructuresDeDonnees.lean\#L192}

\subsection{Arbres binaires stricts}

\begin{definition}
Un arbre binaire strict : chaque nœud interne a exactement deux enfants. On le
décrit par un type à part, ce qui rend la contrainte impossible à violer plutôt que
seulement vérifiée.
\end{definition}
\source{https://github.com/Commutator-IO/learn-lean/blob/main/courses/02-lycee/11-informatique-nsi/StructuresDeDonnees.lean\#L214}{StructuresDeDonnees.lean\#L214}

\begin{definition}
Le nombre de feuilles.
\end{definition}
\source{https://github.com/Commutator-IO/learn-lean/blob/main/courses/02-lycee/11-informatique-nsi/StructuresDeDonnees.lean\#L219}{StructuresDeDonnees.lean\#L219}

\begin{definition}
Le nombre de nœuds internes.
\end{definition}
\source{https://github.com/Commutator-IO/learn-lean/blob/main/courses/02-lycee/11-informatique-nsi/StructuresDeDonnees.lean\#L224}{StructuresDeDonnees.lean\#L224}

\begin{theoreme}
Dans un arbre binaire strict, il y a exactement une feuille de plus que de nœuds
internes. On peut le lire ainsi : chaque nœud interne créé remplace une feuille par
deux, donc ajoute une feuille et un interne ; le compte part de une feuille et zéro
interne.
\end{theoreme}
\begin{proof}
Par récurrence. Une feuille compte une feuille et zéro nœud interne.

Un nœud de sous-arbres \(g\) et \(d\) compte \(f(g) + f(d)\) feuilles et
\(1 + i(g) + i(d)\) internes ; l'hypothèse de récurrence donne \(f(g) = i(g) + 1\) et
\(f(d) = i(d) + 1\), d'où
\[ f = i(g) + i(d) + 2 = \bigl(1 + i(g) + i(d)\bigr) + 1 = i + 1 . \]

On peut le lire ainsi : chaque nœud interne créé remplace une feuille par deux, donc
ajoute une feuille et un interne ; le compte part d'une feuille et zéro interne, et
l'écart de un ne bouge plus.
\end{proof}
\source{https://github.com/Commutator-IO/learn-lean/blob/main/courses/02-lycee/11-informatique-nsi/StructuresDeDonnees.lean\#L232}{StructuresDeDonnees.lean\#L232}

\subsection{Arbres binaires de recherche}

\begin{definition}
Le parcours infixe : le sous-arbre gauche, la racine, le sous-arbre droit.
\end{definition}
\source{https://github.com/Commutator-IO/learn-lean/blob/main/courses/02-lycee/11-informatique-nsi/StructuresDeDonnees.lean\#L242}{StructuresDeDonnees.lean\#L242}

\begin{definition}
La propriété d'arbre binaire de recherche : à chaque nœud, tout ce qui est à
gauche est plus petit que la racine, tout ce qui est à droite est plus grand, et les
deux sous-arbres sont eux-mêmes des arbres binaires de recherche.
\end{definition}
\source{https://github.com/Commutator-IO/learn-lean/blob/main/courses/02-lycee/11-informatique-nsi/StructuresDeDonnees.lean\#L249}{StructuresDeDonnees.lean\#L249}

\begin{theoreme}
Le parcours infixe d'un arbre binaire de recherche rend les clés triées. C'est ce
qui justifie la structure : on obtient le tri sans trier, par le seul fait d'avoir
rangé les clés selon la règle.
\end{theoreme}
\begin{proof}
Par récurrence sur l'arbre. Le parcours de l'arbre vide est la liste vide.

Pour un nœud de racine \(x\), le parcours est la concaténation du parcours gauche, de
\(x\), et du parcours droit. Or une concaténation est triée dès que chaque morceau
l'est et que tout élément du premier est inférieur à tout élément du second.

Les deux parcours sont triés par hypothèse de récurrence. La liste formée de \(x\) suivi
du parcours droit est triée, puisque \(x\) minore ce parcours par définition de l'arbre
binaire de recherche. Enfin, pour \(u\) dans le parcours gauche et \(v\) dans le
parcours droit, \(u \le x \le v\).

C'est ce qui justifie la structure : on obtient le tri sans trier, par le seul fait
d'avoir rangé les clés selon la règle.
\end{proof}
\source{https://github.com/Commutator-IO/learn-lean/blob/main/courses/02-lycee/11-informatique-nsi/StructuresDeDonnees.lean\#L257}{StructuresDeDonnees.lean\#L257}

\begin{figure}[h]
\centering
% Construction (repère mathématique, y vers le haut) :
%   racine 8 en (0, 3) ; 3 en (-2, 2), 10 en (2, 2) ;
%   1 en (-3, 1), 6 en (-1, 1), 14 en (3, 1).
%   Axe du parcours en y = -0.4, chaque clé projetée sur sa propre abscisse.
% SVG : x = 140 + 38x, y = 30 + 38(3 - y).
\begin{tikzpicture}[scale=0.9, every node/.style={font=\small}]
  \tikzstyle{cle}=[draw, circle, inner sep=1.5pt, minimum size=7mm]
  \node[cle] (n8)  at (0, 3)  {$8$};
  \node[cle] (n3)  at (-2, 2) {$3$};
  \node[cle] (n10) at (2, 2)  {$10$};
  \node[cle] (n1)  at (-3, 1) {$1$};
  \node[cle] (n6)  at (-1, 1) {$6$};
  \node[cle] (n14) at (3, 1)  {$14$};
  \draw (n8) -- (n3); \draw (n8) -- (n10);
  \draw (n3) -- (n1); \draw (n3) -- (n6); \draw (n10) -- (n14);
  \draw[dotted] (-3, 0.6) -- (-3, -0.4);
  \draw[dotted] (-2, 1.6) -- (-2, -0.4);
  \draw[dotted] (-1, 0.6) -- (-1, -0.4);
  \draw[dotted] (0, 2.6) -- (0, -0.4);
  \draw[dotted] (2, 1.6) -- (2, -0.4);
  \draw[dotted] (3, 0.6) -- (3, -0.4);
  \draw (-3.6, -0.4) -- (3.6, -0.4);
  \foreach \x/\v in {-3/1, -2/3, -1/6, 0/8, 2/10, 3/14}
    \node[below] at (\x, -0.5) {$\v$};
\end{tikzpicture}

\caption{Un arbre binaire de recherche. Les clés y sont rangées de sorte que le parcours infixe \textemdash{} gauche, racine, droite \textemdash{} les lise de gauche à droite, donc dans l'ordre croissant.}
\end{figure}

\begin{definition}
La recherche dans un arbre binaire de recherche : on compare à la racine et l'on
descend d'un seul côté.
\end{definition}
\source{https://github.com/Commutator-IO/learn-lean/blob/main/courses/02-lycee/11-informatique-nsi/StructuresDeDonnees.lean\#L274}{StructuresDeDonnees.lean\#L274}

\begin{definition}
Le nombre de comparaisons que fait cette recherche.
\end{definition}
\source{https://github.com/Commutator-IO/learn-lean/blob/main/courses/02-lycee/11-informatique-nsi/StructuresDeDonnees.lean\#L279}{StructuresDeDonnees.lean\#L279}

\begin{theoreme}
La recherche dans un arbre binaire de recherche coûte au plus la hauteur de
l'arbre : on ne descend que d'un niveau par comparaison, et l'on ne remonte jamais.
D'où l'intérêt d'un arbre équilibré, dont la hauteur est logarithmique.
\end{theoreme}
\begin{proof}
Par récurrence sur l'arbre. Dans l'arbre vide, aucune comparaison.

Sur un nœud, la recherche compare une fois, puis ou bien s'arrête, ou bien descend dans
un seul des deux sous-arbres. Le coût vaut donc \(1\), ou \(1\) plus le coût dans un
sous-arbre — majoré par hypothèse de récurrence par la hauteur de celui-ci, elle-même
majorée par le maximum des deux hauteurs. Dans tous les cas, le total est au plus
\(1 + \max(h(g), h(d))\), qui est la hauteur de l'arbre.

On ne descend que d'un niveau par comparaison, et l'on ne remonte jamais. D'où l'intérêt
d'un arbre équilibré, dont la hauteur est logarithmique en le nombre de clés.
\end{proof}
\source{https://github.com/Commutator-IO/learn-lean/blob/main/courses/02-lycee/11-informatique-nsi/StructuresDeDonnees.lean\#L287}{StructuresDeDonnees.lean\#L287}

\subsection{Parcours en profondeur d'un graphe}

\begin{definition}
Un graphe est donné par la liste des voisins de chaque sommet. Un sommet \texttt{v} est
accessible depuis \texttt{u} s'il existe une suite d'arêtes qui y mène \textemdash{} y compris la suite
vide, qui rend tout sommet accessible depuis lui-même.
\end{definition}
\source{https://github.com/Commutator-IO/learn-lean/blob/main/courses/02-lycee/11-informatique-nsi/StructuresDeDonnees.lean\#L304}{StructuresDeDonnees.lean\#L304}

\begin{definition}
Le parcours en profondeur, avec une pile explicite et la liste des sommets déjà
vus. Le premier argument est un carburant : il borne le nombre de tours, ce qui évite
d'avoir à démontrer la terminaison avant d'avoir démontré quoi que ce soit d'autre.
Un vrai programme s'arrête parce que la liste des vus ne fait que croître.
\end{definition}
\source{https://github.com/Commutator-IO/learn-lean/blob/main/courses/02-lycee/11-informatique-nsi/StructuresDeDonnees.lean\#L311}{StructuresDeDonnees.lean\#L311}

\begin{theoreme}
Le parcours en profondeur ne visite que des sommets accessibles : il n'invente
rien. La démonstration est une récurrence sur le carburant, avec pour invariant que
tout ce qui est dans la pile et tout ce qui est déjà vu est accessible.
\end{theoreme}
\begin{proof}
Par récurrence sur le carburant, avec pour invariant que tout ce qui est dans la pile et
tout ce qui est déjà vu est accessible depuis le départ.

À carburant nul, le parcours rend la liste des vus, dont l'invariant dit qu'ils sont
accessibles ; si la pile est vide, de même.

Sinon on examine le sommet \(s\) de la pile. S'il est déjà vu, on passe au suivant sans
rien changer. Sinon on le marque vu et l'on empile ses voisins : \(s\) est accessible par
l'invariant, donc ses voisins le sont aussi, d'un pas de plus. L'invariant est conservé,
et l'hypothèse de récurrence conclut.

Le parcours n'invente rien. C'est la moitié facile du théorème.
\end{proof}
\source{https://github.com/Commutator-IO/learn-lean/blob/main/courses/02-lycee/11-informatique-nsi/StructuresDeDonnees.lean\#L322}{StructuresDeDonnees.lean\#L322}

\begin{theoreme}
Réciproquement, tout sommet accessible finit par être visité, pourvu qu'on laisse
au parcours assez de carburant.

Énoncé admis. Sa démonstration demande un invariant nettement plus fin que celui de
la moitié précédente : il faut suivre, à chaque tour, l'ensemble des sommets qui sont
soit déjà vus, soit encore dans la pile, et montrer qu'aucun sommet accessible ne peut
sortir de cet ensemble \textemdash{} puis borner le carburant nécessaire. C'est la moitié
difficile du théorème, et elle n'est pas faite ici.
\end{theoreme}
\admis
\source{https://github.com/Commutator-IO/learn-lean/blob/main/courses/02-lycee/11-informatique-nsi/StructuresDeDonnees.lean\#L359}{StructuresDeDonnees.lean\#L359}

\subsection{Plus courts chemins et algorithme de Dijkstra}

\begin{definition}
Le poids d'un chemin, parcouru depuis un sommet de départ : la somme des poids des
arêtes empruntées. Le chemin est donné par la liste des sommets qu'on visite après le
départ.
\end{definition}
\source{https://github.com/Commutator-IO/learn-lean/blob/main/courses/02-lycee/11-informatique-nsi/StructuresDeDonnees.lean\#L369}{StructuresDeDonnees.lean\#L369}

\begin{theoreme}
Avec des poids positifs, allonger un chemin ne le raccourcit jamais.

C'est exactement ce qui autorise l'algorithme de Dijkstra à fixer définitivement la
distance du sommet le plus proche sans regarder la suite du graphe : aucun détour ne
pourra faire mieux. Avec des poids négatifs, l'énoncé tombe, et l'algorithme aussi.
\end{theoreme}
\begin{proof}
Par récurrence sur le chemin, le sommet de départ restant libre.

Un chemin vide ne coûte rien. Sinon le chemin s'écrit \(v :: r\) et son prolongement
commence par la même arête \(u \to v\) ; l'hypothèse de récurrence compare alors les
restes depuis \(v\).

C'est exactement ce qui autorise l'algorithme de Dijkstra à fixer définitivement la
distance du sommet le plus proche sans regarder la suite du graphe : aucun détour ne
pourra faire mieux. Avec des poids négatifs, l'énoncé tombe, et l'algorithme aussi.
\end{proof}
\source{https://github.com/Commutator-IO/learn-lean/blob/main/courses/02-lycee/11-informatique-nsi/StructuresDeDonnees.lean\#L378}{StructuresDeDonnees.lean\#L378}

\begin{theoreme}
Correction de l'algorithme de Dijkstra pour des poids positifs.

L'algorithme fixe, pour chaque sommet, une distance \texttt{fixe v}. Les deux hypothèses sont
ce qu'il garantit à la fin : la distance de la source est nulle, et aucune arête ne
permet d'améliorer une distance fixée \textemdash{} c'est ce que signifie \og{} l'algorithme a
relâché toutes les arêtes \fg{}. On en tire ce qui est demandé : la distance fixée est
bien la plus petite possible, quel que soit le chemin qu'on considère.

La démonstration remonte le chemin arête par arête depuis la source. Ce que cet
énoncé ne formalise pas : que l'algorithme, avec sa file de priorité, établisse bien
ces deux garanties \textemdash{} c'est la partie programmation, et elle est ici portée par les
hypothèses plutôt que passée sous silence.
\end{theoreme}
\begin{proof}
On démontre d'abord, par récurrence sur le chemin, une inégalité plus générale : pour
tout chemin allant d'un sommet \(u\) à un sommet \(w\),
\[ \mathrm{fixe}(w) \le \mathrm{fixe}(u) + \text{poids du chemin depuis } u . \]

Un chemin vide n'aboutit nulle part : il n'y a rien à démontrer. Si le chemin est réduit
à une arête \(u \to z\), alors \(w = z\) et l'inégalité est exactement l'hypothèse de
relâchement. Sinon le chemin commence par l'arête \(u \to z\) et se poursuit ;
l'hypothèse de récurrence donne \(\mathrm{fixe}(w) \le \mathrm{fixe}(z) + p\), où \(p\)
est le poids du reste, et le relâchement donne
\(\mathrm{fixe}(z) \le \mathrm{fixe}(u) + \mathrm{poids}(u, z)\) ; on additionne.

L'énoncé s'obtient en prenant pour \(u\) la source, dont la distance fixée est nulle.

Les deux hypothèses sont ce que l'algorithme garantit à la fin de son exécution : la
source est à distance nulle, et aucune arête ne permet plus d'améliorer une distance.
Ce que cette démonstration ne formalise pas, c'est que l'algorithme, avec sa file de
priorité, les établisse effectivement.
\end{proof}
\source{https://github.com/Commutator-IO/learn-lean/blob/main/courses/02-lycee/11-informatique-nsi/StructuresDeDonnees.lean\#L398}{StructuresDeDonnees.lean\#L398}

\vfill
\noindent\rule{0.35\linewidth}{0.4pt}\par
\noindent{\footnotesize Leonardo de Moura et Sebastian Ullrich,
\emph{The Lean~4 Theorem Prover and Programming Language},
\emph{Automated Deduction — CADE~28}, Springer, 2021, p.~625--635,
\href{https://doi.org/10.1007/978-3-030-79876-5_37}{doi:10.1007/978-3-030-79876-5\_37}.\par}

\end{document}
